%% shared.tex -- common preamble for the Linguistics of Asian Languages decks.
%%
%% \input this immediately after \documentclass{beamer}.  It fixes the theme,
%% the fonts for every script the course touches, the colour conventions and
%% the language shorthands in ONE place, so that a font fix made here reaches
%% every deck.
%%
%% Deliberately NOT here, because the decks genuinely differ:
%%   * the gb4e flavour -- gb4e, mygb4e and langsci-gb4e are not
%%     interchangeable, so each deck loads the one its examples are written
%%     for, straight after %% shared.tex -- common preamble for the Linguistics of Asian Languages decks.
%%
%% \input this immediately after \documentclass{beamer}.  It fixes the theme,
%% the fonts for every script the course touches, the colour conventions and
%% the language shorthands in ONE place, so that a font fix made here reaches
%% every deck.
%%
%% Deliberately NOT here, because the decks genuinely differ:
%%   * the gb4e flavour -- gb4e, mygb4e and langsci-gb4e are not
%%     interchangeable, so each deck loads the one its examples are written
%%     for, straight after %% shared.tex -- common preamble for the Linguistics of Asian Languages decks.
%%
%% \input this immediately after \documentclass{beamer}.  It fixes the theme,
%% the fonts for every script the course touches, the colour conventions and
%% the language shorthands in ONE place, so that a font fix made here reaches
%% every deck.
%%
%% Deliberately NOT here, because the decks genuinely differ:
%%   * the gb4e flavour -- gb4e, mygb4e and langsci-gb4e are not
%%     interchangeable, so each deck loads the one its examples are written
%%     for, straight after %% shared.tex -- common preamble for the Linguistics of Asian Languages decks.
%%
%% \input this immediately after \documentclass{beamer}.  It fixes the theme,
%% the fonts for every script the course touches, the colour conventions and
%% the language shorthands in ONE place, so that a font fix made here reaches
%% every deck.
%%
%% Deliberately NOT here, because the decks genuinely differ:
%%   * the gb4e flavour -- gb4e, mygb4e and langsci-gb4e are not
%%     interchangeable, so each deck loads the one its examples are written
%%     for, straight after \input{shared}
%%   * \title, and \author where it is not just Francis Bond
%%   * anything a deck deliberately overrides afterwards -- sounds.tex, for
%%     instance, sets Andika as the body font because it is a phonetics deck
%%
%% Everything here can be overridden by the deck: \input{shared} first, then
%% deviate, and put a comment saying why.

%% ---- fonts -------------------------------------------------------------
%%
%% Two mechanisms, deliberately, because the decks use both:
%%
%%   xeCJK      picks the CJK font automatically from the character, so Han,
%%              kana and Hangul just work in running text.
%%   polyglossia gives \textchinese{}, \textjapanese{}, \textkorean{},
%%              \textyue{}, \textthai{} and \textarabic{} for when a specific
%%              regional face is wanted -- the Japanese and Traditional faces
%%              draw many shared characters differently from the Simplified
%%              one, and Arabic needs RTL.
%%
%% Han has no italic.  Map the italic shape onto the upright one for every
%% CJK face, or XeTeX fake-slants it whenever CJK lands inside \textit --
%% which it constantly does, because \mtcitestyle italicises.

\usepackage{fontspec}
\usepackage{xeCJK}
\usepackage{graphicx}
\usepackage{tikz}

% Body font.  Latin Modern Sans, beamer's default, is the wrong font for this
% course: it has no small caps (so \textsc{sg} in a gloss came out as
% lowercase "sg") and it is missing 33 of the characters these decks use --
% most of the IPA, the modifier letters and the super/subscripts -- which
% XeTeX then drops silently.
%
% Libertinus is the maintained successor to Linux Libertine/Biolinum.  Sans
% becomes the body font (beamer sets \sffamily), Serif is there for gloss
% lines, and it brings a real maths font with it.  It lives in TeX Live, so
% kpathsea finds it on any machine without a fontconfig install.
% Set the faces directly rather than via libertinus-otf, which also wants a
% mono font (AnonymousPro) that is not installed here.
% Ligatures={TeX,NoCommon}: keep ``---'' and friends, but switch OFF the fi/fl
% output ligatures, so that text copied out of the PDF comes back as plain
% ASCII.  Those ligatures are how the deck sources picked up "Pacific" and
% "geneDcally" in the first place.
\usefonttheme{professionalfonts}
\setmainfont{LibertinusSerif}[
  Extension		= .otf,
  UprightFont		= *-Regular,
  ItalicFont		= *-Italic,
  BoldFont		= *-Bold,
  BoldItalicFont	= *-BoldItalic,
  Ligatures		= {TeX,NoCommon}]
% Scale 0.95: Libertinus has a much larger x-height than Latin Modern, so at
% the nominal size it overflows frames that used to fit.  0.95 restores every
% deck to the page count it had before the switch.
\setsansfont{LibertinusSans}[
  Scale			= 0.95,
  Extension		= .otf,
  UprightFont		= *-Regular,
  ItalicFont		= *-Italic,
  BoldFont		= *-Bold,
  % Libertinus Sans ships no bold italic, so slant the bold rather than let
  % LaTeX silently substitute an upright bold.
  BoldItalicFont	= *-Bold,
  BoldItalicFeatures	= {FakeSlant = 0.2},
  Ligatures		= {TeX,NoCommon}]
\setmonofont{LibertinusMono-Regular}[Extension = .otf]

% XeTeX always loads face 0 of a TrueType Collection, and
% NotoSansCJK-Regular.ttc packs JP=0, KR=1, SC=2, TC=3, HK=4.  So asking for
% "Noto Sans CJK SC" -- as every deck did -- silently handed back the
% JAPANESE face, and all the Simplified Chinese in this course was being set
% with Japanese glyph shapes (骨, 直, 关 and many more differ).  FontIndex
% picks the right face, and the bare filename resolves without a path.
%
% Han has no italic.  Point ItalicFont at the same face so that CJK landing
% inside \textit -- which happens constantly, because \mtcitestyle
% italicises -- stays upright instead of drawing an "undefined shape".
\newcommand{\CJKface}[2]{%
  \newCJKfontfamily#1{NotoSansCJK-Regular.ttc}[%
    FontIndex = #2, ItalicFont = NotoSansCJK-Regular.ttc,
    ItalicFeatures = {FontIndex = #2}]}

\IfFontExistsTF{NotoSansCJK-Regular.ttc}{%
  \setCJKmainfont{NotoSansCJK-Regular.ttc}[%
    FontIndex = 2, ItalicFont = NotoSansCJK-Regular.ttc,
    ItalicFeatures = {FontIndex = 2}]
  \CJKface\chinesefontsf{2}	% Simplified
  \CJKface\yuefontsf{3}	% Traditional, used for Cantonese
  \CJKface\japanesefontsf{0}
  \CJKface\koreanfontsf{1}
}{%
  % No collection: fall back to the family names.  Everything still builds,
  % but the regional faces will all be whichever one fontconfig hands over.
  \setCJKmainfont{Noto Sans CJK SC}[ItalicFont = Noto Sans CJK SC]
  \newCJKfontfamily\chinesefontsf{Noto Sans CJK SC}[ItalicFont = Noto Sans CJK SC]
  \newCJKfontfamily\yuefontsf{Noto Sans CJK TC}[ItalicFont = Noto Sans CJK TC]
  \newCJKfontfamily\japanesefontsf{Noto Sans CJK JP}[ItalicFont = Noto Sans CJK JP]
  \newCJKfontfamily\koreanfontsf{Noto Sans CJK KR}[ItalicFont = Noto Sans CJK KR]
}
\setCJKsansfont{NotoSansCJK-Regular.ttc}[%
  FontIndex = 2, ItalicFont = NotoSansCJK-Regular.ttc,
  ItalicFeatures = {FontIndex = 2}]
% The mono faces live in the same collection: Mono JP=5, KR=6, SC=7, TC=8.
\setCJKmonofont{NotoSansCJK-Regular.ttc}[FontIndex = 7]

% The bopomofo tone marks sit in Spacing Modifier Letters, not the Bopomofo
% block, so xeCJK would leave them to Latin Modern -- which has no ˊ or ˋ.
\xeCJKDeclareCharClass{CJK}{"02C7, "02C9, "02CA, "02CB, "02D9}

% The reverse problem: xeCJK counts the curly quotes as CJK punctuation, so
% ' ' " " came out full-width in the CJK font, with CJK spacing around them
% and a line-break opportunity inside words like "doesn't".  In this course
% they are always English punctuation.
\xeCJKDeclareCharClass{Default}{"2018, "2019, "201C, "201D}

% IPA needs no font of its own any more: Biolinum covers IPA Extensions,
% Phonetic Extensions and the Spacing Modifier Letters.  The decks had grown
% four different answers to this (Andika, Noto Serif, Libertine, and nothing
% at all), so \ipa, \IPA and \silipa are all kept as names but now resolve
% to the body font.
\newfontfamily\ipafont{LibertinusSans}[
  Extension = .otf, UprightFont = *-Regular, ItalicFont = *-Italic,
  BoldFont = *-Bold, Ligatures = {TeX,NoCommon}]
\newcommand{\ipa}[1]{{\ipafont #1}}
\newcommand{\IPA}[1]{{\ipafont #1}}
\let\silipa\ipafont

% Used for the second line of gb4e glosses, where tone numbers and IPA turn
% up and Latin Modern runs out of superscripts.
\newfontfamily\Libertine{LibertinusSerif}[
  Extension = .otf, UprightFont = *-Regular, ItalicFont = *-Italic,
  BoldFont = *-Bold, BoldItalicFont = *-BoldItalic,
  Ligatures = {TeX,NoCommon}]

\usetheme{Madrid}
\usecolortheme{crane}

\usepackage{polyglossia}
\setdefaultlanguage{english}
\setotherlanguage{thai}
\newfontfamily\thaifont[Script=Thai]{Noto Sans Thai}
\newfontfamily\thaifontsf[Script=Thai]{Noto Sans Thai}
\setotherlanguage{arabic}
\newfontfamily\arabicfont[Script=Arabic]{Noto Naskh Arabic}
\newfontfamily\arabicfontsf[Script=Arabic]{Noto Naskh Arabic}

% polyglossia has no gloss file for these, so define the switches directly
% against the CJK families declared above.
\newcommand{\textchinese}[1]{{\chinesefontsf #1}}
\newcommand{\textjapanese}[1]{{\japanesefontsf #1}}
\newcommand{\textkorean}[1]{{\koreanfontsf #1}}
\newcommand{\textyue}[1]{{\yuefontsf #1}}

% Latin Modern has no mathematical angle brackets, which the writing systems
% deck uses for graphemes.
\usepackage{newunicodechar}
\newunicodechar{⟨}{‹}
\newunicodechar{⟩}{›}

%% ---- colour and emphasis ------------------------------------------------
\usepackage{xcolor}
\newcommand{\txx}[1]{\textcolor{blue}{\textbf{#1}}}
\let\emp\txx
\let\zz\txx
\newcommand{\mtplain}[1]{\textcolor{teal}{#1}}
\newcommand{\con}[1]{\textsc{#1}}

\usepackage[normalem]{ulem}
\newcommand{\ul}[1]{\uline{#1}}
\newcommand{\cll}[1]{\uline{#1}}

%% ---- linguistic examples ------------------------------------------------
%%
%% langsci-gb4e everywhere.  Plain gb4e makes _ and ^ active characters,
%% which breaks hyperref -- that is the one and only reason mygb4e existed, a
%% local copy of gb4e with exactly that hack commented out.  The Language
%% Science Press fork drops it, is maintained, and gives the \exfont /
%% \glossfont / \transfont hooks, so the local fork is no longer needed.
\usepackage{langsci-gb4e}

% Second gloss line -- transcription and IPA -- in the serif face.
% writing.tex and verbal-arts.tex override this afterwards.
\renewcommand{\eachwordtwo}{\Libertine}

%% ---- language shorthands ------------------------------------------------
%%
%% Two families, and the distinction matters:
%%
%%   ISO-code macros  \cmn{}, \jpn{} ... take a ROMANISATION and set it in
%%                    teal italic, optionally with a gloss: \cmn[dog]{gou3}.
%%   \zhs{} and friends take NATIVE SCRIPT and set it in teal upright in the
%%                    right regional face.  Script is never italicised.

\usepackage[c,e,f,j,k]{mtg2e}
\renewcommand{\mtcitestyle}[1]{\textcolor{teal}{\textit{#1}}}

% romanisations (mtg2e already gives \chn, \eng, \fra, \jpn, \kjpn and \kor)
\newcommand{\cmn}{\mtciteform}
\newcommand{\zh}{\mtciteform}
\newcommand{\yue}{\mtciteform}
\newcommand{\ko}{\mtciteform}
\newcommand{\vie}{\mtciteform}
\newcommand{\vi}{\mtciteform}
\newcommand{\tha}{\mtciteform}
\newcommand{\thaig}{\mtciteform}
\newcommand{\lao}{\mtciteform}
\newcommand{\ind}{\mtciteform}
\newcommand{\msa}{\mtciteform}
\newcommand{\zsm}{\mtciteform}
\newcommand{\ar}{\mtciteform}
\newcommand{\lex}[1]{\textbf{\mtcitestyle{#1}}}

% native script
\newcommand{\zhs}[1]{\mtplain{\textchinese{#1}}}
\newcommand{\zht}[1]{\mtplain{\textyue{#1}}}
\newcommand{\jps}[1]{\mtplain{\textjapanese{#1}}}
\newcommand{\kos}[1]{\mtplain{\textkorean{#1}}}
\newcommand{\ths}[1]{\mtplain{\textthai{#1}}}

%% ---- structure -----------------------------------------------------------
%% Sections are signposting, so give each one a page of its own showing where
%% we have got to, with the rest of the deck dimmed.  \overview puts the same
%% list up front; call it straight after \frame{\titlepage}.
\AtBeginSection[]{%
  \begin{frame}{Outline}
    \tableofcontents[currentsection]
  \end{frame}}

\newcommand{\overview}{%
  \begin{frame}{Overview}
    \tableofcontents
  \end{frame}}

%% ---- references ----------------------------------------------------------
\usepackage[round]{natbib}

%% ---- defaults ------------------------------------------------------------
%% The year lives here so the rollover is a one-line change.
\author[Francis Bond]{Francis Bond}
\date{2026}

%%   * \title, and \author where it is not just Francis Bond
%%   * anything a deck deliberately overrides afterwards -- sounds.tex, for
%%     instance, sets Andika as the body font because it is a phonetics deck
%%
%% Everything here can be overridden by the deck: %% shared.tex -- common preamble for the Linguistics of Asian Languages decks.
%%
%% \input this immediately after \documentclass{beamer}.  It fixes the theme,
%% the fonts for every script the course touches, the colour conventions and
%% the language shorthands in ONE place, so that a font fix made here reaches
%% every deck.
%%
%% Deliberately NOT here, because the decks genuinely differ:
%%   * the gb4e flavour -- gb4e, mygb4e and langsci-gb4e are not
%%     interchangeable, so each deck loads the one its examples are written
%%     for, straight after \input{shared}
%%   * \title, and \author where it is not just Francis Bond
%%   * anything a deck deliberately overrides afterwards -- sounds.tex, for
%%     instance, sets Andika as the body font because it is a phonetics deck
%%
%% Everything here can be overridden by the deck: \input{shared} first, then
%% deviate, and put a comment saying why.

%% ---- fonts -------------------------------------------------------------
%%
%% Two mechanisms, deliberately, because the decks use both:
%%
%%   xeCJK      picks the CJK font automatically from the character, so Han,
%%              kana and Hangul just work in running text.
%%   polyglossia gives \textchinese{}, \textjapanese{}, \textkorean{},
%%              \textyue{}, \textthai{} and \textarabic{} for when a specific
%%              regional face is wanted -- the Japanese and Traditional faces
%%              draw many shared characters differently from the Simplified
%%              one, and Arabic needs RTL.
%%
%% Han has no italic.  Map the italic shape onto the upright one for every
%% CJK face, or XeTeX fake-slants it whenever CJK lands inside \textit --
%% which it constantly does, because \mtcitestyle italicises.

\usepackage{fontspec}
\usepackage{xeCJK}
\usepackage{graphicx}
\usepackage{tikz}

% Body font.  Latin Modern Sans, beamer's default, is the wrong font for this
% course: it has no small caps (so \textsc{sg} in a gloss came out as
% lowercase "sg") and it is missing 33 of the characters these decks use --
% most of the IPA, the modifier letters and the super/subscripts -- which
% XeTeX then drops silently.
%
% Libertinus is the maintained successor to Linux Libertine/Biolinum.  Sans
% becomes the body font (beamer sets \sffamily), Serif is there for gloss
% lines, and it brings a real maths font with it.  It lives in TeX Live, so
% kpathsea finds it on any machine without a fontconfig install.
% Set the faces directly rather than via libertinus-otf, which also wants a
% mono font (AnonymousPro) that is not installed here.
% Ligatures={TeX,NoCommon}: keep ``---'' and friends, but switch OFF the fi/fl
% output ligatures, so that text copied out of the PDF comes back as plain
% ASCII.  Those ligatures are how the deck sources picked up "Pacific" and
% "geneDcally" in the first place.
\usefonttheme{professionalfonts}
\setmainfont{LibertinusSerif}[
  Extension		= .otf,
  UprightFont		= *-Regular,
  ItalicFont		= *-Italic,
  BoldFont		= *-Bold,
  BoldItalicFont	= *-BoldItalic,
  Ligatures		= {TeX,NoCommon}]
% Scale 0.95: Libertinus has a much larger x-height than Latin Modern, so at
% the nominal size it overflows frames that used to fit.  0.95 restores every
% deck to the page count it had before the switch.
\setsansfont{LibertinusSans}[
  Scale			= 0.95,
  Extension		= .otf,
  UprightFont		= *-Regular,
  ItalicFont		= *-Italic,
  BoldFont		= *-Bold,
  % Libertinus Sans ships no bold italic, so slant the bold rather than let
  % LaTeX silently substitute an upright bold.
  BoldItalicFont	= *-Bold,
  BoldItalicFeatures	= {FakeSlant = 0.2},
  Ligatures		= {TeX,NoCommon}]
\setmonofont{LibertinusMono-Regular}[Extension = .otf]

% XeTeX always loads face 0 of a TrueType Collection, and
% NotoSansCJK-Regular.ttc packs JP=0, KR=1, SC=2, TC=3, HK=4.  So asking for
% "Noto Sans CJK SC" -- as every deck did -- silently handed back the
% JAPANESE face, and all the Simplified Chinese in this course was being set
% with Japanese glyph shapes (骨, 直, 关 and many more differ).  FontIndex
% picks the right face, and the bare filename resolves without a path.
%
% Han has no italic.  Point ItalicFont at the same face so that CJK landing
% inside \textit -- which happens constantly, because \mtcitestyle
% italicises -- stays upright instead of drawing an "undefined shape".
\newcommand{\CJKface}[2]{%
  \newCJKfontfamily#1{NotoSansCJK-Regular.ttc}[%
    FontIndex = #2, ItalicFont = NotoSansCJK-Regular.ttc,
    ItalicFeatures = {FontIndex = #2}]}

\IfFontExistsTF{NotoSansCJK-Regular.ttc}{%
  \setCJKmainfont{NotoSansCJK-Regular.ttc}[%
    FontIndex = 2, ItalicFont = NotoSansCJK-Regular.ttc,
    ItalicFeatures = {FontIndex = 2}]
  \CJKface\chinesefontsf{2}	% Simplified
  \CJKface\yuefontsf{3}	% Traditional, used for Cantonese
  \CJKface\japanesefontsf{0}
  \CJKface\koreanfontsf{1}
}{%
  % No collection: fall back to the family names.  Everything still builds,
  % but the regional faces will all be whichever one fontconfig hands over.
  \setCJKmainfont{Noto Sans CJK SC}[ItalicFont = Noto Sans CJK SC]
  \newCJKfontfamily\chinesefontsf{Noto Sans CJK SC}[ItalicFont = Noto Sans CJK SC]
  \newCJKfontfamily\yuefontsf{Noto Sans CJK TC}[ItalicFont = Noto Sans CJK TC]
  \newCJKfontfamily\japanesefontsf{Noto Sans CJK JP}[ItalicFont = Noto Sans CJK JP]
  \newCJKfontfamily\koreanfontsf{Noto Sans CJK KR}[ItalicFont = Noto Sans CJK KR]
}
\setCJKsansfont{NotoSansCJK-Regular.ttc}[%
  FontIndex = 2, ItalicFont = NotoSansCJK-Regular.ttc,
  ItalicFeatures = {FontIndex = 2}]
% The mono faces live in the same collection: Mono JP=5, KR=6, SC=7, TC=8.
\setCJKmonofont{NotoSansCJK-Regular.ttc}[FontIndex = 7]

% The bopomofo tone marks sit in Spacing Modifier Letters, not the Bopomofo
% block, so xeCJK would leave them to Latin Modern -- which has no ˊ or ˋ.
\xeCJKDeclareCharClass{CJK}{"02C7, "02C9, "02CA, "02CB, "02D9}

% The reverse problem: xeCJK counts the curly quotes as CJK punctuation, so
% ' ' " " came out full-width in the CJK font, with CJK spacing around them
% and a line-break opportunity inside words like "doesn't".  In this course
% they are always English punctuation.
\xeCJKDeclareCharClass{Default}{"2018, "2019, "201C, "201D}

% IPA needs no font of its own any more: Biolinum covers IPA Extensions,
% Phonetic Extensions and the Spacing Modifier Letters.  The decks had grown
% four different answers to this (Andika, Noto Serif, Libertine, and nothing
% at all), so \ipa, \IPA and \silipa are all kept as names but now resolve
% to the body font.
\newfontfamily\ipafont{LibertinusSans}[
  Extension = .otf, UprightFont = *-Regular, ItalicFont = *-Italic,
  BoldFont = *-Bold, Ligatures = {TeX,NoCommon}]
\newcommand{\ipa}[1]{{\ipafont #1}}
\newcommand{\IPA}[1]{{\ipafont #1}}
\let\silipa\ipafont

% Used for the second line of gb4e glosses, where tone numbers and IPA turn
% up and Latin Modern runs out of superscripts.
\newfontfamily\Libertine{LibertinusSerif}[
  Extension = .otf, UprightFont = *-Regular, ItalicFont = *-Italic,
  BoldFont = *-Bold, BoldItalicFont = *-BoldItalic,
  Ligatures = {TeX,NoCommon}]

\usetheme{Madrid}
\usecolortheme{crane}

\usepackage{polyglossia}
\setdefaultlanguage{english}
\setotherlanguage{thai}
\newfontfamily\thaifont[Script=Thai]{Noto Sans Thai}
\newfontfamily\thaifontsf[Script=Thai]{Noto Sans Thai}
\setotherlanguage{arabic}
\newfontfamily\arabicfont[Script=Arabic]{Noto Naskh Arabic}
\newfontfamily\arabicfontsf[Script=Arabic]{Noto Naskh Arabic}

% polyglossia has no gloss file for these, so define the switches directly
% against the CJK families declared above.
\newcommand{\textchinese}[1]{{\chinesefontsf #1}}
\newcommand{\textjapanese}[1]{{\japanesefontsf #1}}
\newcommand{\textkorean}[1]{{\koreanfontsf #1}}
\newcommand{\textyue}[1]{{\yuefontsf #1}}

% Latin Modern has no mathematical angle brackets, which the writing systems
% deck uses for graphemes.
\usepackage{newunicodechar}
\newunicodechar{⟨}{‹}
\newunicodechar{⟩}{›}

%% ---- colour and emphasis ------------------------------------------------
\usepackage{xcolor}
\newcommand{\txx}[1]{\textcolor{blue}{\textbf{#1}}}
\let\emp\txx
\let\zz\txx
\newcommand{\mtplain}[1]{\textcolor{teal}{#1}}
\newcommand{\con}[1]{\textsc{#1}}

\usepackage[normalem]{ulem}
\newcommand{\ul}[1]{\uline{#1}}
\newcommand{\cll}[1]{\uline{#1}}

%% ---- linguistic examples ------------------------------------------------
%%
%% langsci-gb4e everywhere.  Plain gb4e makes _ and ^ active characters,
%% which breaks hyperref -- that is the one and only reason mygb4e existed, a
%% local copy of gb4e with exactly that hack commented out.  The Language
%% Science Press fork drops it, is maintained, and gives the \exfont /
%% \glossfont / \transfont hooks, so the local fork is no longer needed.
\usepackage{langsci-gb4e}

% Second gloss line -- transcription and IPA -- in the serif face.
% writing.tex and verbal-arts.tex override this afterwards.
\renewcommand{\eachwordtwo}{\Libertine}

%% ---- language shorthands ------------------------------------------------
%%
%% Two families, and the distinction matters:
%%
%%   ISO-code macros  \cmn{}, \jpn{} ... take a ROMANISATION and set it in
%%                    teal italic, optionally with a gloss: \cmn[dog]{gou3}.
%%   \zhs{} and friends take NATIVE SCRIPT and set it in teal upright in the
%%                    right regional face.  Script is never italicised.

\usepackage[c,e,f,j,k]{mtg2e}
\renewcommand{\mtcitestyle}[1]{\textcolor{teal}{\textit{#1}}}

% romanisations (mtg2e already gives \chn, \eng, \fra, \jpn, \kjpn and \kor)
\newcommand{\cmn}{\mtciteform}
\newcommand{\zh}{\mtciteform}
\newcommand{\yue}{\mtciteform}
\newcommand{\ko}{\mtciteform}
\newcommand{\vie}{\mtciteform}
\newcommand{\vi}{\mtciteform}
\newcommand{\tha}{\mtciteform}
\newcommand{\thaig}{\mtciteform}
\newcommand{\lao}{\mtciteform}
\newcommand{\ind}{\mtciteform}
\newcommand{\msa}{\mtciteform}
\newcommand{\zsm}{\mtciteform}
\newcommand{\ar}{\mtciteform}
\newcommand{\lex}[1]{\textbf{\mtcitestyle{#1}}}

% native script
\newcommand{\zhs}[1]{\mtplain{\textchinese{#1}}}
\newcommand{\zht}[1]{\mtplain{\textyue{#1}}}
\newcommand{\jps}[1]{\mtplain{\textjapanese{#1}}}
\newcommand{\kos}[1]{\mtplain{\textkorean{#1}}}
\newcommand{\ths}[1]{\mtplain{\textthai{#1}}}

%% ---- structure -----------------------------------------------------------
%% Sections are signposting, so give each one a page of its own showing where
%% we have got to, with the rest of the deck dimmed.  \overview puts the same
%% list up front; call it straight after \frame{\titlepage}.
\AtBeginSection[]{%
  \begin{frame}{Outline}
    \tableofcontents[currentsection]
  \end{frame}}

\newcommand{\overview}{%
  \begin{frame}{Overview}
    \tableofcontents
  \end{frame}}

%% ---- references ----------------------------------------------------------
\usepackage[round]{natbib}

%% ---- defaults ------------------------------------------------------------
%% The year lives here so the rollover is a one-line change.
\author[Francis Bond]{Francis Bond}
\date{2026}
 first, then
%% deviate, and put a comment saying why.

%% ---- fonts -------------------------------------------------------------
%%
%% Two mechanisms, deliberately, because the decks use both:
%%
%%   xeCJK      picks the CJK font automatically from the character, so Han,
%%              kana and Hangul just work in running text.
%%   polyglossia gives \textchinese{}, \textjapanese{}, \textkorean{},
%%              \textyue{}, \textthai{} and \textarabic{} for when a specific
%%              regional face is wanted -- the Japanese and Traditional faces
%%              draw many shared characters differently from the Simplified
%%              one, and Arabic needs RTL.
%%
%% Han has no italic.  Map the italic shape onto the upright one for every
%% CJK face, or XeTeX fake-slants it whenever CJK lands inside \textit --
%% which it constantly does, because \mtcitestyle italicises.

\usepackage{fontspec}
\usepackage{xeCJK}
\usepackage{graphicx}
\usepackage{tikz}

% Body font.  Latin Modern Sans, beamer's default, is the wrong font for this
% course: it has no small caps (so \textsc{sg} in a gloss came out as
% lowercase "sg") and it is missing 33 of the characters these decks use --
% most of the IPA, the modifier letters and the super/subscripts -- which
% XeTeX then drops silently.
%
% Libertinus is the maintained successor to Linux Libertine/Biolinum.  Sans
% becomes the body font (beamer sets \sffamily), Serif is there for gloss
% lines, and it brings a real maths font with it.  It lives in TeX Live, so
% kpathsea finds it on any machine without a fontconfig install.
% Set the faces directly rather than via libertinus-otf, which also wants a
% mono font (AnonymousPro) that is not installed here.
% Ligatures={TeX,NoCommon}: keep ``---'' and friends, but switch OFF the fi/fl
% output ligatures, so that text copied out of the PDF comes back as plain
% ASCII.  Those ligatures are how the deck sources picked up "Pacific" and
% "geneDcally" in the first place.
\usefonttheme{professionalfonts}
\setmainfont{LibertinusSerif}[
  Extension		= .otf,
  UprightFont		= *-Regular,
  ItalicFont		= *-Italic,
  BoldFont		= *-Bold,
  BoldItalicFont	= *-BoldItalic,
  Ligatures		= {TeX,NoCommon}]
% Scale 0.95: Libertinus has a much larger x-height than Latin Modern, so at
% the nominal size it overflows frames that used to fit.  0.95 restores every
% deck to the page count it had before the switch.
\setsansfont{LibertinusSans}[
  Scale			= 0.95,
  Extension		= .otf,
  UprightFont		= *-Regular,
  ItalicFont		= *-Italic,
  BoldFont		= *-Bold,
  % Libertinus Sans ships no bold italic, so slant the bold rather than let
  % LaTeX silently substitute an upright bold.
  BoldItalicFont	= *-Bold,
  BoldItalicFeatures	= {FakeSlant = 0.2},
  Ligatures		= {TeX,NoCommon}]
\setmonofont{LibertinusMono-Regular}[Extension = .otf]

% XeTeX always loads face 0 of a TrueType Collection, and
% NotoSansCJK-Regular.ttc packs JP=0, KR=1, SC=2, TC=3, HK=4.  So asking for
% "Noto Sans CJK SC" -- as every deck did -- silently handed back the
% JAPANESE face, and all the Simplified Chinese in this course was being set
% with Japanese glyph shapes (骨, 直, 关 and many more differ).  FontIndex
% picks the right face, and the bare filename resolves without a path.
%
% Han has no italic.  Point ItalicFont at the same face so that CJK landing
% inside \textit -- which happens constantly, because \mtcitestyle
% italicises -- stays upright instead of drawing an "undefined shape".
\newcommand{\CJKface}[2]{%
  \newCJKfontfamily#1{NotoSansCJK-Regular.ttc}[%
    FontIndex = #2, ItalicFont = NotoSansCJK-Regular.ttc,
    ItalicFeatures = {FontIndex = #2}]}

\IfFontExistsTF{NotoSansCJK-Regular.ttc}{%
  \setCJKmainfont{NotoSansCJK-Regular.ttc}[%
    FontIndex = 2, ItalicFont = NotoSansCJK-Regular.ttc,
    ItalicFeatures = {FontIndex = 2}]
  \CJKface\chinesefontsf{2}	% Simplified
  \CJKface\yuefontsf{3}	% Traditional, used for Cantonese
  \CJKface\japanesefontsf{0}
  \CJKface\koreanfontsf{1}
}{%
  % No collection: fall back to the family names.  Everything still builds,
  % but the regional faces will all be whichever one fontconfig hands over.
  \setCJKmainfont{Noto Sans CJK SC}[ItalicFont = Noto Sans CJK SC]
  \newCJKfontfamily\chinesefontsf{Noto Sans CJK SC}[ItalicFont = Noto Sans CJK SC]
  \newCJKfontfamily\yuefontsf{Noto Sans CJK TC}[ItalicFont = Noto Sans CJK TC]
  \newCJKfontfamily\japanesefontsf{Noto Sans CJK JP}[ItalicFont = Noto Sans CJK JP]
  \newCJKfontfamily\koreanfontsf{Noto Sans CJK KR}[ItalicFont = Noto Sans CJK KR]
}
\setCJKsansfont{NotoSansCJK-Regular.ttc}[%
  FontIndex = 2, ItalicFont = NotoSansCJK-Regular.ttc,
  ItalicFeatures = {FontIndex = 2}]
% The mono faces live in the same collection: Mono JP=5, KR=6, SC=7, TC=8.
\setCJKmonofont{NotoSansCJK-Regular.ttc}[FontIndex = 7]

% The bopomofo tone marks sit in Spacing Modifier Letters, not the Bopomofo
% block, so xeCJK would leave them to Latin Modern -- which has no ˊ or ˋ.
\xeCJKDeclareCharClass{CJK}{"02C7, "02C9, "02CA, "02CB, "02D9}

% The reverse problem: xeCJK counts the curly quotes as CJK punctuation, so
% ' ' " " came out full-width in the CJK font, with CJK spacing around them
% and a line-break opportunity inside words like "doesn't".  In this course
% they are always English punctuation.
\xeCJKDeclareCharClass{Default}{"2018, "2019, "201C, "201D}

% IPA needs no font of its own any more: Biolinum covers IPA Extensions,
% Phonetic Extensions and the Spacing Modifier Letters.  The decks had grown
% four different answers to this (Andika, Noto Serif, Libertine, and nothing
% at all), so \ipa, \IPA and \silipa are all kept as names but now resolve
% to the body font.
\newfontfamily\ipafont{LibertinusSans}[
  Extension = .otf, UprightFont = *-Regular, ItalicFont = *-Italic,
  BoldFont = *-Bold, Ligatures = {TeX,NoCommon}]
\newcommand{\ipa}[1]{{\ipafont #1}}
\newcommand{\IPA}[1]{{\ipafont #1}}
\let\silipa\ipafont

% Used for the second line of gb4e glosses, where tone numbers and IPA turn
% up and Latin Modern runs out of superscripts.
\newfontfamily\Libertine{LibertinusSerif}[
  Extension = .otf, UprightFont = *-Regular, ItalicFont = *-Italic,
  BoldFont = *-Bold, BoldItalicFont = *-BoldItalic,
  Ligatures = {TeX,NoCommon}]

\usetheme{Madrid}
\usecolortheme{crane}

\usepackage{polyglossia}
\setdefaultlanguage{english}
\setotherlanguage{thai}
\newfontfamily\thaifont[Script=Thai]{Noto Sans Thai}
\newfontfamily\thaifontsf[Script=Thai]{Noto Sans Thai}
\setotherlanguage{arabic}
\newfontfamily\arabicfont[Script=Arabic]{Noto Naskh Arabic}
\newfontfamily\arabicfontsf[Script=Arabic]{Noto Naskh Arabic}

% polyglossia has no gloss file for these, so define the switches directly
% against the CJK families declared above.
\newcommand{\textchinese}[1]{{\chinesefontsf #1}}
\newcommand{\textjapanese}[1]{{\japanesefontsf #1}}
\newcommand{\textkorean}[1]{{\koreanfontsf #1}}
\newcommand{\textyue}[1]{{\yuefontsf #1}}

% Latin Modern has no mathematical angle brackets, which the writing systems
% deck uses for graphemes.
\usepackage{newunicodechar}
\newunicodechar{⟨}{‹}
\newunicodechar{⟩}{›}

%% ---- colour and emphasis ------------------------------------------------
\usepackage{xcolor}
\newcommand{\txx}[1]{\textcolor{blue}{\textbf{#1}}}
\let\emp\txx
\let\zz\txx
\newcommand{\mtplain}[1]{\textcolor{teal}{#1}}
\newcommand{\con}[1]{\textsc{#1}}

\usepackage[normalem]{ulem}
\newcommand{\ul}[1]{\uline{#1}}
\newcommand{\cll}[1]{\uline{#1}}

%% ---- linguistic examples ------------------------------------------------
%%
%% langsci-gb4e everywhere.  Plain gb4e makes _ and ^ active characters,
%% which breaks hyperref -- that is the one and only reason mygb4e existed, a
%% local copy of gb4e with exactly that hack commented out.  The Language
%% Science Press fork drops it, is maintained, and gives the \exfont /
%% \glossfont / \transfont hooks, so the local fork is no longer needed.
\usepackage{langsci-gb4e}

% Second gloss line -- transcription and IPA -- in the serif face.
% writing.tex and verbal-arts.tex override this afterwards.
\renewcommand{\eachwordtwo}{\Libertine}

%% ---- language shorthands ------------------------------------------------
%%
%% Two families, and the distinction matters:
%%
%%   ISO-code macros  \cmn{}, \jpn{} ... take a ROMANISATION and set it in
%%                    teal italic, optionally with a gloss: \cmn[dog]{gou3}.
%%   \zhs{} and friends take NATIVE SCRIPT and set it in teal upright in the
%%                    right regional face.  Script is never italicised.

\usepackage[c,e,f,j,k]{mtg2e}
\renewcommand{\mtcitestyle}[1]{\textcolor{teal}{\textit{#1}}}

% romanisations (mtg2e already gives \chn, \eng, \fra, \jpn, \kjpn and \kor)
\newcommand{\cmn}{\mtciteform}
\newcommand{\zh}{\mtciteform}
\newcommand{\yue}{\mtciteform}
\newcommand{\ko}{\mtciteform}
\newcommand{\vie}{\mtciteform}
\newcommand{\vi}{\mtciteform}
\newcommand{\tha}{\mtciteform}
\newcommand{\thaig}{\mtciteform}
\newcommand{\lao}{\mtciteform}
\newcommand{\ind}{\mtciteform}
\newcommand{\msa}{\mtciteform}
\newcommand{\zsm}{\mtciteform}
\newcommand{\ar}{\mtciteform}
\newcommand{\lex}[1]{\textbf{\mtcitestyle{#1}}}

% native script
\newcommand{\zhs}[1]{\mtplain{\textchinese{#1}}}
\newcommand{\zht}[1]{\mtplain{\textyue{#1}}}
\newcommand{\jps}[1]{\mtplain{\textjapanese{#1}}}
\newcommand{\kos}[1]{\mtplain{\textkorean{#1}}}
\newcommand{\ths}[1]{\mtplain{\textthai{#1}}}

%% ---- structure -----------------------------------------------------------
%% Sections are signposting, so give each one a page of its own showing where
%% we have got to, with the rest of the deck dimmed.  \overview puts the same
%% list up front; call it straight after \frame{\titlepage}.
\AtBeginSection[]{%
  \begin{frame}{Outline}
    \tableofcontents[currentsection]
  \end{frame}}

\newcommand{\overview}{%
  \begin{frame}{Overview}
    \tableofcontents
  \end{frame}}

%% ---- references ----------------------------------------------------------
\usepackage[round]{natbib}

%% ---- defaults ------------------------------------------------------------
%% The year lives here so the rollover is a one-line change.
\author[Francis Bond]{Francis Bond}
\date{2026}

%%   * \title, and \author where it is not just Francis Bond
%%   * anything a deck deliberately overrides afterwards -- sounds.tex, for
%%     instance, sets Andika as the body font because it is a phonetics deck
%%
%% Everything here can be overridden by the deck: %% shared.tex -- common preamble for the Linguistics of Asian Languages decks.
%%
%% \input this immediately after \documentclass{beamer}.  It fixes the theme,
%% the fonts for every script the course touches, the colour conventions and
%% the language shorthands in ONE place, so that a font fix made here reaches
%% every deck.
%%
%% Deliberately NOT here, because the decks genuinely differ:
%%   * the gb4e flavour -- gb4e, mygb4e and langsci-gb4e are not
%%     interchangeable, so each deck loads the one its examples are written
%%     for, straight after %% shared.tex -- common preamble for the Linguistics of Asian Languages decks.
%%
%% \input this immediately after \documentclass{beamer}.  It fixes the theme,
%% the fonts for every script the course touches, the colour conventions and
%% the language shorthands in ONE place, so that a font fix made here reaches
%% every deck.
%%
%% Deliberately NOT here, because the decks genuinely differ:
%%   * the gb4e flavour -- gb4e, mygb4e and langsci-gb4e are not
%%     interchangeable, so each deck loads the one its examples are written
%%     for, straight after \input{shared}
%%   * \title, and \author where it is not just Francis Bond
%%   * anything a deck deliberately overrides afterwards -- sounds.tex, for
%%     instance, sets Andika as the body font because it is a phonetics deck
%%
%% Everything here can be overridden by the deck: \input{shared} first, then
%% deviate, and put a comment saying why.

%% ---- fonts -------------------------------------------------------------
%%
%% Two mechanisms, deliberately, because the decks use both:
%%
%%   xeCJK      picks the CJK font automatically from the character, so Han,
%%              kana and Hangul just work in running text.
%%   polyglossia gives \textchinese{}, \textjapanese{}, \textkorean{},
%%              \textyue{}, \textthai{} and \textarabic{} for when a specific
%%              regional face is wanted -- the Japanese and Traditional faces
%%              draw many shared characters differently from the Simplified
%%              one, and Arabic needs RTL.
%%
%% Han has no italic.  Map the italic shape onto the upright one for every
%% CJK face, or XeTeX fake-slants it whenever CJK lands inside \textit --
%% which it constantly does, because \mtcitestyle italicises.

\usepackage{fontspec}
\usepackage{xeCJK}
\usepackage{graphicx}
\usepackage{tikz}

% Body font.  Latin Modern Sans, beamer's default, is the wrong font for this
% course: it has no small caps (so \textsc{sg} in a gloss came out as
% lowercase "sg") and it is missing 33 of the characters these decks use --
% most of the IPA, the modifier letters and the super/subscripts -- which
% XeTeX then drops silently.
%
% Libertinus is the maintained successor to Linux Libertine/Biolinum.  Sans
% becomes the body font (beamer sets \sffamily), Serif is there for gloss
% lines, and it brings a real maths font with it.  It lives in TeX Live, so
% kpathsea finds it on any machine without a fontconfig install.
% Set the faces directly rather than via libertinus-otf, which also wants a
% mono font (AnonymousPro) that is not installed here.
% Ligatures={TeX,NoCommon}: keep ``---'' and friends, but switch OFF the fi/fl
% output ligatures, so that text copied out of the PDF comes back as plain
% ASCII.  Those ligatures are how the deck sources picked up "Pacific" and
% "geneDcally" in the first place.
\usefonttheme{professionalfonts}
\setmainfont{LibertinusSerif}[
  Extension		= .otf,
  UprightFont		= *-Regular,
  ItalicFont		= *-Italic,
  BoldFont		= *-Bold,
  BoldItalicFont	= *-BoldItalic,
  Ligatures		= {TeX,NoCommon}]
% Scale 0.95: Libertinus has a much larger x-height than Latin Modern, so at
% the nominal size it overflows frames that used to fit.  0.95 restores every
% deck to the page count it had before the switch.
\setsansfont{LibertinusSans}[
  Scale			= 0.95,
  Extension		= .otf,
  UprightFont		= *-Regular,
  ItalicFont		= *-Italic,
  BoldFont		= *-Bold,
  % Libertinus Sans ships no bold italic, so slant the bold rather than let
  % LaTeX silently substitute an upright bold.
  BoldItalicFont	= *-Bold,
  BoldItalicFeatures	= {FakeSlant = 0.2},
  Ligatures		= {TeX,NoCommon}]
\setmonofont{LibertinusMono-Regular}[Extension = .otf]

% XeTeX always loads face 0 of a TrueType Collection, and
% NotoSansCJK-Regular.ttc packs JP=0, KR=1, SC=2, TC=3, HK=4.  So asking for
% "Noto Sans CJK SC" -- as every deck did -- silently handed back the
% JAPANESE face, and all the Simplified Chinese in this course was being set
% with Japanese glyph shapes (骨, 直, 关 and many more differ).  FontIndex
% picks the right face, and the bare filename resolves without a path.
%
% Han has no italic.  Point ItalicFont at the same face so that CJK landing
% inside \textit -- which happens constantly, because \mtcitestyle
% italicises -- stays upright instead of drawing an "undefined shape".
\newcommand{\CJKface}[2]{%
  \newCJKfontfamily#1{NotoSansCJK-Regular.ttc}[%
    FontIndex = #2, ItalicFont = NotoSansCJK-Regular.ttc,
    ItalicFeatures = {FontIndex = #2}]}

\IfFontExistsTF{NotoSansCJK-Regular.ttc}{%
  \setCJKmainfont{NotoSansCJK-Regular.ttc}[%
    FontIndex = 2, ItalicFont = NotoSansCJK-Regular.ttc,
    ItalicFeatures = {FontIndex = 2}]
  \CJKface\chinesefontsf{2}	% Simplified
  \CJKface\yuefontsf{3}	% Traditional, used for Cantonese
  \CJKface\japanesefontsf{0}
  \CJKface\koreanfontsf{1}
}{%
  % No collection: fall back to the family names.  Everything still builds,
  % but the regional faces will all be whichever one fontconfig hands over.
  \setCJKmainfont{Noto Sans CJK SC}[ItalicFont = Noto Sans CJK SC]
  \newCJKfontfamily\chinesefontsf{Noto Sans CJK SC}[ItalicFont = Noto Sans CJK SC]
  \newCJKfontfamily\yuefontsf{Noto Sans CJK TC}[ItalicFont = Noto Sans CJK TC]
  \newCJKfontfamily\japanesefontsf{Noto Sans CJK JP}[ItalicFont = Noto Sans CJK JP]
  \newCJKfontfamily\koreanfontsf{Noto Sans CJK KR}[ItalicFont = Noto Sans CJK KR]
}
\setCJKsansfont{NotoSansCJK-Regular.ttc}[%
  FontIndex = 2, ItalicFont = NotoSansCJK-Regular.ttc,
  ItalicFeatures = {FontIndex = 2}]
% The mono faces live in the same collection: Mono JP=5, KR=6, SC=7, TC=8.
\setCJKmonofont{NotoSansCJK-Regular.ttc}[FontIndex = 7]

% The bopomofo tone marks sit in Spacing Modifier Letters, not the Bopomofo
% block, so xeCJK would leave them to Latin Modern -- which has no ˊ or ˋ.
\xeCJKDeclareCharClass{CJK}{"02C7, "02C9, "02CA, "02CB, "02D9}

% The reverse problem: xeCJK counts the curly quotes as CJK punctuation, so
% ' ' " " came out full-width in the CJK font, with CJK spacing around them
% and a line-break opportunity inside words like "doesn't".  In this course
% they are always English punctuation.
\xeCJKDeclareCharClass{Default}{"2018, "2019, "201C, "201D}

% IPA needs no font of its own any more: Biolinum covers IPA Extensions,
% Phonetic Extensions and the Spacing Modifier Letters.  The decks had grown
% four different answers to this (Andika, Noto Serif, Libertine, and nothing
% at all), so \ipa, \IPA and \silipa are all kept as names but now resolve
% to the body font.
\newfontfamily\ipafont{LibertinusSans}[
  Extension = .otf, UprightFont = *-Regular, ItalicFont = *-Italic,
  BoldFont = *-Bold, Ligatures = {TeX,NoCommon}]
\newcommand{\ipa}[1]{{\ipafont #1}}
\newcommand{\IPA}[1]{{\ipafont #1}}
\let\silipa\ipafont

% Used for the second line of gb4e glosses, where tone numbers and IPA turn
% up and Latin Modern runs out of superscripts.
\newfontfamily\Libertine{LibertinusSerif}[
  Extension = .otf, UprightFont = *-Regular, ItalicFont = *-Italic,
  BoldFont = *-Bold, BoldItalicFont = *-BoldItalic,
  Ligatures = {TeX,NoCommon}]

\usetheme{Madrid}
\usecolortheme{crane}

\usepackage{polyglossia}
\setdefaultlanguage{english}
\setotherlanguage{thai}
\newfontfamily\thaifont[Script=Thai]{Noto Sans Thai}
\newfontfamily\thaifontsf[Script=Thai]{Noto Sans Thai}
\setotherlanguage{arabic}
\newfontfamily\arabicfont[Script=Arabic]{Noto Naskh Arabic}
\newfontfamily\arabicfontsf[Script=Arabic]{Noto Naskh Arabic}

% polyglossia has no gloss file for these, so define the switches directly
% against the CJK families declared above.
\newcommand{\textchinese}[1]{{\chinesefontsf #1}}
\newcommand{\textjapanese}[1]{{\japanesefontsf #1}}
\newcommand{\textkorean}[1]{{\koreanfontsf #1}}
\newcommand{\textyue}[1]{{\yuefontsf #1}}

% Latin Modern has no mathematical angle brackets, which the writing systems
% deck uses for graphemes.
\usepackage{newunicodechar}
\newunicodechar{⟨}{‹}
\newunicodechar{⟩}{›}

%% ---- colour and emphasis ------------------------------------------------
\usepackage{xcolor}
\newcommand{\txx}[1]{\textcolor{blue}{\textbf{#1}}}
\let\emp\txx
\let\zz\txx
\newcommand{\mtplain}[1]{\textcolor{teal}{#1}}
\newcommand{\con}[1]{\textsc{#1}}

\usepackage[normalem]{ulem}
\newcommand{\ul}[1]{\uline{#1}}
\newcommand{\cll}[1]{\uline{#1}}

%% ---- linguistic examples ------------------------------------------------
%%
%% langsci-gb4e everywhere.  Plain gb4e makes _ and ^ active characters,
%% which breaks hyperref -- that is the one and only reason mygb4e existed, a
%% local copy of gb4e with exactly that hack commented out.  The Language
%% Science Press fork drops it, is maintained, and gives the \exfont /
%% \glossfont / \transfont hooks, so the local fork is no longer needed.
\usepackage{langsci-gb4e}

% Second gloss line -- transcription and IPA -- in the serif face.
% writing.tex and verbal-arts.tex override this afterwards.
\renewcommand{\eachwordtwo}{\Libertine}

%% ---- language shorthands ------------------------------------------------
%%
%% Two families, and the distinction matters:
%%
%%   ISO-code macros  \cmn{}, \jpn{} ... take a ROMANISATION and set it in
%%                    teal italic, optionally with a gloss: \cmn[dog]{gou3}.
%%   \zhs{} and friends take NATIVE SCRIPT and set it in teal upright in the
%%                    right regional face.  Script is never italicised.

\usepackage[c,e,f,j,k]{mtg2e}
\renewcommand{\mtcitestyle}[1]{\textcolor{teal}{\textit{#1}}}

% romanisations (mtg2e already gives \chn, \eng, \fra, \jpn, \kjpn and \kor)
\newcommand{\cmn}{\mtciteform}
\newcommand{\zh}{\mtciteform}
\newcommand{\yue}{\mtciteform}
\newcommand{\ko}{\mtciteform}
\newcommand{\vie}{\mtciteform}
\newcommand{\vi}{\mtciteform}
\newcommand{\tha}{\mtciteform}
\newcommand{\thaig}{\mtciteform}
\newcommand{\lao}{\mtciteform}
\newcommand{\ind}{\mtciteform}
\newcommand{\msa}{\mtciteform}
\newcommand{\zsm}{\mtciteform}
\newcommand{\ar}{\mtciteform}
\newcommand{\lex}[1]{\textbf{\mtcitestyle{#1}}}

% native script
\newcommand{\zhs}[1]{\mtplain{\textchinese{#1}}}
\newcommand{\zht}[1]{\mtplain{\textyue{#1}}}
\newcommand{\jps}[1]{\mtplain{\textjapanese{#1}}}
\newcommand{\kos}[1]{\mtplain{\textkorean{#1}}}
\newcommand{\ths}[1]{\mtplain{\textthai{#1}}}

%% ---- structure -----------------------------------------------------------
%% Sections are signposting, so give each one a page of its own showing where
%% we have got to, with the rest of the deck dimmed.  \overview puts the same
%% list up front; call it straight after \frame{\titlepage}.
\AtBeginSection[]{%
  \begin{frame}{Outline}
    \tableofcontents[currentsection]
  \end{frame}}

\newcommand{\overview}{%
  \begin{frame}{Overview}
    \tableofcontents
  \end{frame}}

%% ---- references ----------------------------------------------------------
\usepackage[round]{natbib}

%% ---- defaults ------------------------------------------------------------
%% The year lives here so the rollover is a one-line change.
\author[Francis Bond]{Francis Bond}
\date{2026}

%%   * \title, and \author where it is not just Francis Bond
%%   * anything a deck deliberately overrides afterwards -- sounds.tex, for
%%     instance, sets Andika as the body font because it is a phonetics deck
%%
%% Everything here can be overridden by the deck: %% shared.tex -- common preamble for the Linguistics of Asian Languages decks.
%%
%% \input this immediately after \documentclass{beamer}.  It fixes the theme,
%% the fonts for every script the course touches, the colour conventions and
%% the language shorthands in ONE place, so that a font fix made here reaches
%% every deck.
%%
%% Deliberately NOT here, because the decks genuinely differ:
%%   * the gb4e flavour -- gb4e, mygb4e and langsci-gb4e are not
%%     interchangeable, so each deck loads the one its examples are written
%%     for, straight after \input{shared}
%%   * \title, and \author where it is not just Francis Bond
%%   * anything a deck deliberately overrides afterwards -- sounds.tex, for
%%     instance, sets Andika as the body font because it is a phonetics deck
%%
%% Everything here can be overridden by the deck: \input{shared} first, then
%% deviate, and put a comment saying why.

%% ---- fonts -------------------------------------------------------------
%%
%% Two mechanisms, deliberately, because the decks use both:
%%
%%   xeCJK      picks the CJK font automatically from the character, so Han,
%%              kana and Hangul just work in running text.
%%   polyglossia gives \textchinese{}, \textjapanese{}, \textkorean{},
%%              \textyue{}, \textthai{} and \textarabic{} for when a specific
%%              regional face is wanted -- the Japanese and Traditional faces
%%              draw many shared characters differently from the Simplified
%%              one, and Arabic needs RTL.
%%
%% Han has no italic.  Map the italic shape onto the upright one for every
%% CJK face, or XeTeX fake-slants it whenever CJK lands inside \textit --
%% which it constantly does, because \mtcitestyle italicises.

\usepackage{fontspec}
\usepackage{xeCJK}
\usepackage{graphicx}
\usepackage{tikz}

% Body font.  Latin Modern Sans, beamer's default, is the wrong font for this
% course: it has no small caps (so \textsc{sg} in a gloss came out as
% lowercase "sg") and it is missing 33 of the characters these decks use --
% most of the IPA, the modifier letters and the super/subscripts -- which
% XeTeX then drops silently.
%
% Libertinus is the maintained successor to Linux Libertine/Biolinum.  Sans
% becomes the body font (beamer sets \sffamily), Serif is there for gloss
% lines, and it brings a real maths font with it.  It lives in TeX Live, so
% kpathsea finds it on any machine without a fontconfig install.
% Set the faces directly rather than via libertinus-otf, which also wants a
% mono font (AnonymousPro) that is not installed here.
% Ligatures={TeX,NoCommon}: keep ``---'' and friends, but switch OFF the fi/fl
% output ligatures, so that text copied out of the PDF comes back as plain
% ASCII.  Those ligatures are how the deck sources picked up "Pacific" and
% "geneDcally" in the first place.
\usefonttheme{professionalfonts}
\setmainfont{LibertinusSerif}[
  Extension		= .otf,
  UprightFont		= *-Regular,
  ItalicFont		= *-Italic,
  BoldFont		= *-Bold,
  BoldItalicFont	= *-BoldItalic,
  Ligatures		= {TeX,NoCommon}]
% Scale 0.95: Libertinus has a much larger x-height than Latin Modern, so at
% the nominal size it overflows frames that used to fit.  0.95 restores every
% deck to the page count it had before the switch.
\setsansfont{LibertinusSans}[
  Scale			= 0.95,
  Extension		= .otf,
  UprightFont		= *-Regular,
  ItalicFont		= *-Italic,
  BoldFont		= *-Bold,
  % Libertinus Sans ships no bold italic, so slant the bold rather than let
  % LaTeX silently substitute an upright bold.
  BoldItalicFont	= *-Bold,
  BoldItalicFeatures	= {FakeSlant = 0.2},
  Ligatures		= {TeX,NoCommon}]
\setmonofont{LibertinusMono-Regular}[Extension = .otf]

% XeTeX always loads face 0 of a TrueType Collection, and
% NotoSansCJK-Regular.ttc packs JP=0, KR=1, SC=2, TC=3, HK=4.  So asking for
% "Noto Sans CJK SC" -- as every deck did -- silently handed back the
% JAPANESE face, and all the Simplified Chinese in this course was being set
% with Japanese glyph shapes (骨, 直, 关 and many more differ).  FontIndex
% picks the right face, and the bare filename resolves without a path.
%
% Han has no italic.  Point ItalicFont at the same face so that CJK landing
% inside \textit -- which happens constantly, because \mtcitestyle
% italicises -- stays upright instead of drawing an "undefined shape".
\newcommand{\CJKface}[2]{%
  \newCJKfontfamily#1{NotoSansCJK-Regular.ttc}[%
    FontIndex = #2, ItalicFont = NotoSansCJK-Regular.ttc,
    ItalicFeatures = {FontIndex = #2}]}

\IfFontExistsTF{NotoSansCJK-Regular.ttc}{%
  \setCJKmainfont{NotoSansCJK-Regular.ttc}[%
    FontIndex = 2, ItalicFont = NotoSansCJK-Regular.ttc,
    ItalicFeatures = {FontIndex = 2}]
  \CJKface\chinesefontsf{2}	% Simplified
  \CJKface\yuefontsf{3}	% Traditional, used for Cantonese
  \CJKface\japanesefontsf{0}
  \CJKface\koreanfontsf{1}
}{%
  % No collection: fall back to the family names.  Everything still builds,
  % but the regional faces will all be whichever one fontconfig hands over.
  \setCJKmainfont{Noto Sans CJK SC}[ItalicFont = Noto Sans CJK SC]
  \newCJKfontfamily\chinesefontsf{Noto Sans CJK SC}[ItalicFont = Noto Sans CJK SC]
  \newCJKfontfamily\yuefontsf{Noto Sans CJK TC}[ItalicFont = Noto Sans CJK TC]
  \newCJKfontfamily\japanesefontsf{Noto Sans CJK JP}[ItalicFont = Noto Sans CJK JP]
  \newCJKfontfamily\koreanfontsf{Noto Sans CJK KR}[ItalicFont = Noto Sans CJK KR]
}
\setCJKsansfont{NotoSansCJK-Regular.ttc}[%
  FontIndex = 2, ItalicFont = NotoSansCJK-Regular.ttc,
  ItalicFeatures = {FontIndex = 2}]
% The mono faces live in the same collection: Mono JP=5, KR=6, SC=7, TC=8.
\setCJKmonofont{NotoSansCJK-Regular.ttc}[FontIndex = 7]

% The bopomofo tone marks sit in Spacing Modifier Letters, not the Bopomofo
% block, so xeCJK would leave them to Latin Modern -- which has no ˊ or ˋ.
\xeCJKDeclareCharClass{CJK}{"02C7, "02C9, "02CA, "02CB, "02D9}

% The reverse problem: xeCJK counts the curly quotes as CJK punctuation, so
% ' ' " " came out full-width in the CJK font, with CJK spacing around them
% and a line-break opportunity inside words like "doesn't".  In this course
% they are always English punctuation.
\xeCJKDeclareCharClass{Default}{"2018, "2019, "201C, "201D}

% IPA needs no font of its own any more: Biolinum covers IPA Extensions,
% Phonetic Extensions and the Spacing Modifier Letters.  The decks had grown
% four different answers to this (Andika, Noto Serif, Libertine, and nothing
% at all), so \ipa, \IPA and \silipa are all kept as names but now resolve
% to the body font.
\newfontfamily\ipafont{LibertinusSans}[
  Extension = .otf, UprightFont = *-Regular, ItalicFont = *-Italic,
  BoldFont = *-Bold, Ligatures = {TeX,NoCommon}]
\newcommand{\ipa}[1]{{\ipafont #1}}
\newcommand{\IPA}[1]{{\ipafont #1}}
\let\silipa\ipafont

% Used for the second line of gb4e glosses, where tone numbers and IPA turn
% up and Latin Modern runs out of superscripts.
\newfontfamily\Libertine{LibertinusSerif}[
  Extension = .otf, UprightFont = *-Regular, ItalicFont = *-Italic,
  BoldFont = *-Bold, BoldItalicFont = *-BoldItalic,
  Ligatures = {TeX,NoCommon}]

\usetheme{Madrid}
\usecolortheme{crane}

\usepackage{polyglossia}
\setdefaultlanguage{english}
\setotherlanguage{thai}
\newfontfamily\thaifont[Script=Thai]{Noto Sans Thai}
\newfontfamily\thaifontsf[Script=Thai]{Noto Sans Thai}
\setotherlanguage{arabic}
\newfontfamily\arabicfont[Script=Arabic]{Noto Naskh Arabic}
\newfontfamily\arabicfontsf[Script=Arabic]{Noto Naskh Arabic}

% polyglossia has no gloss file for these, so define the switches directly
% against the CJK families declared above.
\newcommand{\textchinese}[1]{{\chinesefontsf #1}}
\newcommand{\textjapanese}[1]{{\japanesefontsf #1}}
\newcommand{\textkorean}[1]{{\koreanfontsf #1}}
\newcommand{\textyue}[1]{{\yuefontsf #1}}

% Latin Modern has no mathematical angle brackets, which the writing systems
% deck uses for graphemes.
\usepackage{newunicodechar}
\newunicodechar{⟨}{‹}
\newunicodechar{⟩}{›}

%% ---- colour and emphasis ------------------------------------------------
\usepackage{xcolor}
\newcommand{\txx}[1]{\textcolor{blue}{\textbf{#1}}}
\let\emp\txx
\let\zz\txx
\newcommand{\mtplain}[1]{\textcolor{teal}{#1}}
\newcommand{\con}[1]{\textsc{#1}}

\usepackage[normalem]{ulem}
\newcommand{\ul}[1]{\uline{#1}}
\newcommand{\cll}[1]{\uline{#1}}

%% ---- linguistic examples ------------------------------------------------
%%
%% langsci-gb4e everywhere.  Plain gb4e makes _ and ^ active characters,
%% which breaks hyperref -- that is the one and only reason mygb4e existed, a
%% local copy of gb4e with exactly that hack commented out.  The Language
%% Science Press fork drops it, is maintained, and gives the \exfont /
%% \glossfont / \transfont hooks, so the local fork is no longer needed.
\usepackage{langsci-gb4e}

% Second gloss line -- transcription and IPA -- in the serif face.
% writing.tex and verbal-arts.tex override this afterwards.
\renewcommand{\eachwordtwo}{\Libertine}

%% ---- language shorthands ------------------------------------------------
%%
%% Two families, and the distinction matters:
%%
%%   ISO-code macros  \cmn{}, \jpn{} ... take a ROMANISATION and set it in
%%                    teal italic, optionally with a gloss: \cmn[dog]{gou3}.
%%   \zhs{} and friends take NATIVE SCRIPT and set it in teal upright in the
%%                    right regional face.  Script is never italicised.

\usepackage[c,e,f,j,k]{mtg2e}
\renewcommand{\mtcitestyle}[1]{\textcolor{teal}{\textit{#1}}}

% romanisations (mtg2e already gives \chn, \eng, \fra, \jpn, \kjpn and \kor)
\newcommand{\cmn}{\mtciteform}
\newcommand{\zh}{\mtciteform}
\newcommand{\yue}{\mtciteform}
\newcommand{\ko}{\mtciteform}
\newcommand{\vie}{\mtciteform}
\newcommand{\vi}{\mtciteform}
\newcommand{\tha}{\mtciteform}
\newcommand{\thaig}{\mtciteform}
\newcommand{\lao}{\mtciteform}
\newcommand{\ind}{\mtciteform}
\newcommand{\msa}{\mtciteform}
\newcommand{\zsm}{\mtciteform}
\newcommand{\ar}{\mtciteform}
\newcommand{\lex}[1]{\textbf{\mtcitestyle{#1}}}

% native script
\newcommand{\zhs}[1]{\mtplain{\textchinese{#1}}}
\newcommand{\zht}[1]{\mtplain{\textyue{#1}}}
\newcommand{\jps}[1]{\mtplain{\textjapanese{#1}}}
\newcommand{\kos}[1]{\mtplain{\textkorean{#1}}}
\newcommand{\ths}[1]{\mtplain{\textthai{#1}}}

%% ---- structure -----------------------------------------------------------
%% Sections are signposting, so give each one a page of its own showing where
%% we have got to, with the rest of the deck dimmed.  \overview puts the same
%% list up front; call it straight after \frame{\titlepage}.
\AtBeginSection[]{%
  \begin{frame}{Outline}
    \tableofcontents[currentsection]
  \end{frame}}

\newcommand{\overview}{%
  \begin{frame}{Overview}
    \tableofcontents
  \end{frame}}

%% ---- references ----------------------------------------------------------
\usepackage[round]{natbib}

%% ---- defaults ------------------------------------------------------------
%% The year lives here so the rollover is a one-line change.
\author[Francis Bond]{Francis Bond}
\date{2026}
 first, then
%% deviate, and put a comment saying why.

%% ---- fonts -------------------------------------------------------------
%%
%% Two mechanisms, deliberately, because the decks use both:
%%
%%   xeCJK      picks the CJK font automatically from the character, so Han,
%%              kana and Hangul just work in running text.
%%   polyglossia gives \textchinese{}, \textjapanese{}, \textkorean{},
%%              \textyue{}, \textthai{} and \textarabic{} for when a specific
%%              regional face is wanted -- the Japanese and Traditional faces
%%              draw many shared characters differently from the Simplified
%%              one, and Arabic needs RTL.
%%
%% Han has no italic.  Map the italic shape onto the upright one for every
%% CJK face, or XeTeX fake-slants it whenever CJK lands inside \textit --
%% which it constantly does, because \mtcitestyle italicises.

\usepackage{fontspec}
\usepackage{xeCJK}
\usepackage{graphicx}
\usepackage{tikz}

% Body font.  Latin Modern Sans, beamer's default, is the wrong font for this
% course: it has no small caps (so \textsc{sg} in a gloss came out as
% lowercase "sg") and it is missing 33 of the characters these decks use --
% most of the IPA, the modifier letters and the super/subscripts -- which
% XeTeX then drops silently.
%
% Libertinus is the maintained successor to Linux Libertine/Biolinum.  Sans
% becomes the body font (beamer sets \sffamily), Serif is there for gloss
% lines, and it brings a real maths font with it.  It lives in TeX Live, so
% kpathsea finds it on any machine without a fontconfig install.
% Set the faces directly rather than via libertinus-otf, which also wants a
% mono font (AnonymousPro) that is not installed here.
% Ligatures={TeX,NoCommon}: keep ``---'' and friends, but switch OFF the fi/fl
% output ligatures, so that text copied out of the PDF comes back as plain
% ASCII.  Those ligatures are how the deck sources picked up "Pacific" and
% "geneDcally" in the first place.
\usefonttheme{professionalfonts}
\setmainfont{LibertinusSerif}[
  Extension		= .otf,
  UprightFont		= *-Regular,
  ItalicFont		= *-Italic,
  BoldFont		= *-Bold,
  BoldItalicFont	= *-BoldItalic,
  Ligatures		= {TeX,NoCommon}]
% Scale 0.95: Libertinus has a much larger x-height than Latin Modern, so at
% the nominal size it overflows frames that used to fit.  0.95 restores every
% deck to the page count it had before the switch.
\setsansfont{LibertinusSans}[
  Scale			= 0.95,
  Extension		= .otf,
  UprightFont		= *-Regular,
  ItalicFont		= *-Italic,
  BoldFont		= *-Bold,
  % Libertinus Sans ships no bold italic, so slant the bold rather than let
  % LaTeX silently substitute an upright bold.
  BoldItalicFont	= *-Bold,
  BoldItalicFeatures	= {FakeSlant = 0.2},
  Ligatures		= {TeX,NoCommon}]
\setmonofont{LibertinusMono-Regular}[Extension = .otf]

% XeTeX always loads face 0 of a TrueType Collection, and
% NotoSansCJK-Regular.ttc packs JP=0, KR=1, SC=2, TC=3, HK=4.  So asking for
% "Noto Sans CJK SC" -- as every deck did -- silently handed back the
% JAPANESE face, and all the Simplified Chinese in this course was being set
% with Japanese glyph shapes (骨, 直, 关 and many more differ).  FontIndex
% picks the right face, and the bare filename resolves without a path.
%
% Han has no italic.  Point ItalicFont at the same face so that CJK landing
% inside \textit -- which happens constantly, because \mtcitestyle
% italicises -- stays upright instead of drawing an "undefined shape".
\newcommand{\CJKface}[2]{%
  \newCJKfontfamily#1{NotoSansCJK-Regular.ttc}[%
    FontIndex = #2, ItalicFont = NotoSansCJK-Regular.ttc,
    ItalicFeatures = {FontIndex = #2}]}

\IfFontExistsTF{NotoSansCJK-Regular.ttc}{%
  \setCJKmainfont{NotoSansCJK-Regular.ttc}[%
    FontIndex = 2, ItalicFont = NotoSansCJK-Regular.ttc,
    ItalicFeatures = {FontIndex = 2}]
  \CJKface\chinesefontsf{2}	% Simplified
  \CJKface\yuefontsf{3}	% Traditional, used for Cantonese
  \CJKface\japanesefontsf{0}
  \CJKface\koreanfontsf{1}
}{%
  % No collection: fall back to the family names.  Everything still builds,
  % but the regional faces will all be whichever one fontconfig hands over.
  \setCJKmainfont{Noto Sans CJK SC}[ItalicFont = Noto Sans CJK SC]
  \newCJKfontfamily\chinesefontsf{Noto Sans CJK SC}[ItalicFont = Noto Sans CJK SC]
  \newCJKfontfamily\yuefontsf{Noto Sans CJK TC}[ItalicFont = Noto Sans CJK TC]
  \newCJKfontfamily\japanesefontsf{Noto Sans CJK JP}[ItalicFont = Noto Sans CJK JP]
  \newCJKfontfamily\koreanfontsf{Noto Sans CJK KR}[ItalicFont = Noto Sans CJK KR]
}
\setCJKsansfont{NotoSansCJK-Regular.ttc}[%
  FontIndex = 2, ItalicFont = NotoSansCJK-Regular.ttc,
  ItalicFeatures = {FontIndex = 2}]
% The mono faces live in the same collection: Mono JP=5, KR=6, SC=7, TC=8.
\setCJKmonofont{NotoSansCJK-Regular.ttc}[FontIndex = 7]

% The bopomofo tone marks sit in Spacing Modifier Letters, not the Bopomofo
% block, so xeCJK would leave them to Latin Modern -- which has no ˊ or ˋ.
\xeCJKDeclareCharClass{CJK}{"02C7, "02C9, "02CA, "02CB, "02D9}

% The reverse problem: xeCJK counts the curly quotes as CJK punctuation, so
% ' ' " " came out full-width in the CJK font, with CJK spacing around them
% and a line-break opportunity inside words like "doesn't".  In this course
% they are always English punctuation.
\xeCJKDeclareCharClass{Default}{"2018, "2019, "201C, "201D}

% IPA needs no font of its own any more: Biolinum covers IPA Extensions,
% Phonetic Extensions and the Spacing Modifier Letters.  The decks had grown
% four different answers to this (Andika, Noto Serif, Libertine, and nothing
% at all), so \ipa, \IPA and \silipa are all kept as names but now resolve
% to the body font.
\newfontfamily\ipafont{LibertinusSans}[
  Extension = .otf, UprightFont = *-Regular, ItalicFont = *-Italic,
  BoldFont = *-Bold, Ligatures = {TeX,NoCommon}]
\newcommand{\ipa}[1]{{\ipafont #1}}
\newcommand{\IPA}[1]{{\ipafont #1}}
\let\silipa\ipafont

% Used for the second line of gb4e glosses, where tone numbers and IPA turn
% up and Latin Modern runs out of superscripts.
\newfontfamily\Libertine{LibertinusSerif}[
  Extension = .otf, UprightFont = *-Regular, ItalicFont = *-Italic,
  BoldFont = *-Bold, BoldItalicFont = *-BoldItalic,
  Ligatures = {TeX,NoCommon}]

\usetheme{Madrid}
\usecolortheme{crane}

\usepackage{polyglossia}
\setdefaultlanguage{english}
\setotherlanguage{thai}
\newfontfamily\thaifont[Script=Thai]{Noto Sans Thai}
\newfontfamily\thaifontsf[Script=Thai]{Noto Sans Thai}
\setotherlanguage{arabic}
\newfontfamily\arabicfont[Script=Arabic]{Noto Naskh Arabic}
\newfontfamily\arabicfontsf[Script=Arabic]{Noto Naskh Arabic}

% polyglossia has no gloss file for these, so define the switches directly
% against the CJK families declared above.
\newcommand{\textchinese}[1]{{\chinesefontsf #1}}
\newcommand{\textjapanese}[1]{{\japanesefontsf #1}}
\newcommand{\textkorean}[1]{{\koreanfontsf #1}}
\newcommand{\textyue}[1]{{\yuefontsf #1}}

% Latin Modern has no mathematical angle brackets, which the writing systems
% deck uses for graphemes.
\usepackage{newunicodechar}
\newunicodechar{⟨}{‹}
\newunicodechar{⟩}{›}

%% ---- colour and emphasis ------------------------------------------------
\usepackage{xcolor}
\newcommand{\txx}[1]{\textcolor{blue}{\textbf{#1}}}
\let\emp\txx
\let\zz\txx
\newcommand{\mtplain}[1]{\textcolor{teal}{#1}}
\newcommand{\con}[1]{\textsc{#1}}

\usepackage[normalem]{ulem}
\newcommand{\ul}[1]{\uline{#1}}
\newcommand{\cll}[1]{\uline{#1}}

%% ---- linguistic examples ------------------------------------------------
%%
%% langsci-gb4e everywhere.  Plain gb4e makes _ and ^ active characters,
%% which breaks hyperref -- that is the one and only reason mygb4e existed, a
%% local copy of gb4e with exactly that hack commented out.  The Language
%% Science Press fork drops it, is maintained, and gives the \exfont /
%% \glossfont / \transfont hooks, so the local fork is no longer needed.
\usepackage{langsci-gb4e}

% Second gloss line -- transcription and IPA -- in the serif face.
% writing.tex and verbal-arts.tex override this afterwards.
\renewcommand{\eachwordtwo}{\Libertine}

%% ---- language shorthands ------------------------------------------------
%%
%% Two families, and the distinction matters:
%%
%%   ISO-code macros  \cmn{}, \jpn{} ... take a ROMANISATION and set it in
%%                    teal italic, optionally with a gloss: \cmn[dog]{gou3}.
%%   \zhs{} and friends take NATIVE SCRIPT and set it in teal upright in the
%%                    right regional face.  Script is never italicised.

\usepackage[c,e,f,j,k]{mtg2e}
\renewcommand{\mtcitestyle}[1]{\textcolor{teal}{\textit{#1}}}

% romanisations (mtg2e already gives \chn, \eng, \fra, \jpn, \kjpn and \kor)
\newcommand{\cmn}{\mtciteform}
\newcommand{\zh}{\mtciteform}
\newcommand{\yue}{\mtciteform}
\newcommand{\ko}{\mtciteform}
\newcommand{\vie}{\mtciteform}
\newcommand{\vi}{\mtciteform}
\newcommand{\tha}{\mtciteform}
\newcommand{\thaig}{\mtciteform}
\newcommand{\lao}{\mtciteform}
\newcommand{\ind}{\mtciteform}
\newcommand{\msa}{\mtciteform}
\newcommand{\zsm}{\mtciteform}
\newcommand{\ar}{\mtciteform}
\newcommand{\lex}[1]{\textbf{\mtcitestyle{#1}}}

% native script
\newcommand{\zhs}[1]{\mtplain{\textchinese{#1}}}
\newcommand{\zht}[1]{\mtplain{\textyue{#1}}}
\newcommand{\jps}[1]{\mtplain{\textjapanese{#1}}}
\newcommand{\kos}[1]{\mtplain{\textkorean{#1}}}
\newcommand{\ths}[1]{\mtplain{\textthai{#1}}}

%% ---- structure -----------------------------------------------------------
%% Sections are signposting, so give each one a page of its own showing where
%% we have got to, with the rest of the deck dimmed.  \overview puts the same
%% list up front; call it straight after \frame{\titlepage}.
\AtBeginSection[]{%
  \begin{frame}{Outline}
    \tableofcontents[currentsection]
  \end{frame}}

\newcommand{\overview}{%
  \begin{frame}{Overview}
    \tableofcontents
  \end{frame}}

%% ---- references ----------------------------------------------------------
\usepackage[round]{natbib}

%% ---- defaults ------------------------------------------------------------
%% The year lives here so the rollover is a one-line change.
\author[Francis Bond]{Francis Bond}
\date{2026}
 first, then
%% deviate, and put a comment saying why.

%% ---- fonts -------------------------------------------------------------
%%
%% Two mechanisms, deliberately, because the decks use both:
%%
%%   xeCJK      picks the CJK font automatically from the character, so Han,
%%              kana and Hangul just work in running text.
%%   polyglossia gives \textchinese{}, \textjapanese{}, \textkorean{},
%%              \textyue{}, \textthai{} and \textarabic{} for when a specific
%%              regional face is wanted -- the Japanese and Traditional faces
%%              draw many shared characters differently from the Simplified
%%              one, and Arabic needs RTL.
%%
%% Han has no italic.  Map the italic shape onto the upright one for every
%% CJK face, or XeTeX fake-slants it whenever CJK lands inside \textit --
%% which it constantly does, because \mtcitestyle italicises.

\usepackage{fontspec}
\usepackage{xeCJK}
\usepackage{graphicx}
\usepackage{tikz}

% Body font.  Latin Modern Sans, beamer's default, is the wrong font for this
% course: it has no small caps (so \textsc{sg} in a gloss came out as
% lowercase "sg") and it is missing 33 of the characters these decks use --
% most of the IPA, the modifier letters and the super/subscripts -- which
% XeTeX then drops silently.
%
% Libertinus is the maintained successor to Linux Libertine/Biolinum.  Sans
% becomes the body font (beamer sets \sffamily), Serif is there for gloss
% lines, and it brings a real maths font with it.  It lives in TeX Live, so
% kpathsea finds it on any machine without a fontconfig install.
% Set the faces directly rather than via libertinus-otf, which also wants a
% mono font (AnonymousPro) that is not installed here.
% Ligatures={TeX,NoCommon}: keep ``---'' and friends, but switch OFF the fi/fl
% output ligatures, so that text copied out of the PDF comes back as plain
% ASCII.  Those ligatures are how the deck sources picked up "Pacific" and
% "geneDcally" in the first place.
\usefonttheme{professionalfonts}
\setmainfont{LibertinusSerif}[
  Extension		= .otf,
  UprightFont		= *-Regular,
  ItalicFont		= *-Italic,
  BoldFont		= *-Bold,
  BoldItalicFont	= *-BoldItalic,
  Ligatures		= {TeX,NoCommon}]
% Scale 0.95: Libertinus has a much larger x-height than Latin Modern, so at
% the nominal size it overflows frames that used to fit.  0.95 restores every
% deck to the page count it had before the switch.
\setsansfont{LibertinusSans}[
  Scale			= 0.95,
  Extension		= .otf,
  UprightFont		= *-Regular,
  ItalicFont		= *-Italic,
  BoldFont		= *-Bold,
  % Libertinus Sans ships no bold italic, so slant the bold rather than let
  % LaTeX silently substitute an upright bold.
  BoldItalicFont	= *-Bold,
  BoldItalicFeatures	= {FakeSlant = 0.2},
  Ligatures		= {TeX,NoCommon}]
\setmonofont{LibertinusMono-Regular}[Extension = .otf]

% XeTeX always loads face 0 of a TrueType Collection, and
% NotoSansCJK-Regular.ttc packs JP=0, KR=1, SC=2, TC=3, HK=4.  So asking for
% "Noto Sans CJK SC" -- as every deck did -- silently handed back the
% JAPANESE face, and all the Simplified Chinese in this course was being set
% with Japanese glyph shapes (骨, 直, 关 and many more differ).  FontIndex
% picks the right face, and the bare filename resolves without a path.
%
% Han has no italic.  Point ItalicFont at the same face so that CJK landing
% inside \textit -- which happens constantly, because \mtcitestyle
% italicises -- stays upright instead of drawing an "undefined shape".
\newcommand{\CJKface}[2]{%
  \newCJKfontfamily#1{NotoSansCJK-Regular.ttc}[%
    FontIndex = #2, ItalicFont = NotoSansCJK-Regular.ttc,
    ItalicFeatures = {FontIndex = #2}]}

\IfFontExistsTF{NotoSansCJK-Regular.ttc}{%
  \setCJKmainfont{NotoSansCJK-Regular.ttc}[%
    FontIndex = 2, ItalicFont = NotoSansCJK-Regular.ttc,
    ItalicFeatures = {FontIndex = 2}]
  \CJKface\chinesefontsf{2}	% Simplified
  \CJKface\yuefontsf{3}	% Traditional, used for Cantonese
  \CJKface\japanesefontsf{0}
  \CJKface\koreanfontsf{1}
}{%
  % No collection: fall back to the family names.  Everything still builds,
  % but the regional faces will all be whichever one fontconfig hands over.
  \setCJKmainfont{Noto Sans CJK SC}[ItalicFont = Noto Sans CJK SC]
  \newCJKfontfamily\chinesefontsf{Noto Sans CJK SC}[ItalicFont = Noto Sans CJK SC]
  \newCJKfontfamily\yuefontsf{Noto Sans CJK TC}[ItalicFont = Noto Sans CJK TC]
  \newCJKfontfamily\japanesefontsf{Noto Sans CJK JP}[ItalicFont = Noto Sans CJK JP]
  \newCJKfontfamily\koreanfontsf{Noto Sans CJK KR}[ItalicFont = Noto Sans CJK KR]
}
\setCJKsansfont{NotoSansCJK-Regular.ttc}[%
  FontIndex = 2, ItalicFont = NotoSansCJK-Regular.ttc,
  ItalicFeatures = {FontIndex = 2}]
% The mono faces live in the same collection: Mono JP=5, KR=6, SC=7, TC=8.
\setCJKmonofont{NotoSansCJK-Regular.ttc}[FontIndex = 7]

% The bopomofo tone marks sit in Spacing Modifier Letters, not the Bopomofo
% block, so xeCJK would leave them to Latin Modern -- which has no ˊ or ˋ.
\xeCJKDeclareCharClass{CJK}{"02C7, "02C9, "02CA, "02CB, "02D9}

% The reverse problem: xeCJK counts the curly quotes as CJK punctuation, so
% ' ' " " came out full-width in the CJK font, with CJK spacing around them
% and a line-break opportunity inside words like "doesn't".  In this course
% they are always English punctuation.
\xeCJKDeclareCharClass{Default}{"2018, "2019, "201C, "201D}

% IPA needs no font of its own any more: Biolinum covers IPA Extensions,
% Phonetic Extensions and the Spacing Modifier Letters.  The decks had grown
% four different answers to this (Andika, Noto Serif, Libertine, and nothing
% at all), so \ipa, \IPA and \silipa are all kept as names but now resolve
% to the body font.
\newfontfamily\ipafont{LibertinusSans}[
  Extension = .otf, UprightFont = *-Regular, ItalicFont = *-Italic,
  BoldFont = *-Bold, Ligatures = {TeX,NoCommon}]
\newcommand{\ipa}[1]{{\ipafont #1}}
\newcommand{\IPA}[1]{{\ipafont #1}}
\let\silipa\ipafont

% Used for the second line of gb4e glosses, where tone numbers and IPA turn
% up and Latin Modern runs out of superscripts.
\newfontfamily\Libertine{LibertinusSerif}[
  Extension = .otf, UprightFont = *-Regular, ItalicFont = *-Italic,
  BoldFont = *-Bold, BoldItalicFont = *-BoldItalic,
  Ligatures = {TeX,NoCommon}]

\usetheme{Madrid}
\usecolortheme{crane}

\usepackage{polyglossia}
\setdefaultlanguage{english}
\setotherlanguage{thai}
\newfontfamily\thaifont[Script=Thai]{Noto Sans Thai}
\newfontfamily\thaifontsf[Script=Thai]{Noto Sans Thai}
\setotherlanguage{arabic}
\newfontfamily\arabicfont[Script=Arabic]{Noto Naskh Arabic}
\newfontfamily\arabicfontsf[Script=Arabic]{Noto Naskh Arabic}

% polyglossia has no gloss file for these, so define the switches directly
% against the CJK families declared above.
\newcommand{\textchinese}[1]{{\chinesefontsf #1}}
\newcommand{\textjapanese}[1]{{\japanesefontsf #1}}
\newcommand{\textkorean}[1]{{\koreanfontsf #1}}
\newcommand{\textyue}[1]{{\yuefontsf #1}}

% Latin Modern has no mathematical angle brackets, which the writing systems
% deck uses for graphemes.
\usepackage{newunicodechar}
\newunicodechar{⟨}{‹}
\newunicodechar{⟩}{›}

%% ---- colour and emphasis ------------------------------------------------
\usepackage{xcolor}
\newcommand{\txx}[1]{\textcolor{blue}{\textbf{#1}}}
\let\emp\txx
\let\zz\txx
\newcommand{\mtplain}[1]{\textcolor{teal}{#1}}
\newcommand{\con}[1]{\textsc{#1}}

\usepackage[normalem]{ulem}
\newcommand{\ul}[1]{\uline{#1}}
\newcommand{\cll}[1]{\uline{#1}}

%% ---- linguistic examples ------------------------------------------------
%%
%% langsci-gb4e everywhere.  Plain gb4e makes _ and ^ active characters,
%% which breaks hyperref -- that is the one and only reason mygb4e existed, a
%% local copy of gb4e with exactly that hack commented out.  The Language
%% Science Press fork drops it, is maintained, and gives the \exfont /
%% \glossfont / \transfont hooks, so the local fork is no longer needed.
\usepackage{langsci-gb4e}

% Second gloss line -- transcription and IPA -- in the serif face.
% writing.tex and verbal-arts.tex override this afterwards.
\renewcommand{\eachwordtwo}{\Libertine}

%% ---- language shorthands ------------------------------------------------
%%
%% Two families, and the distinction matters:
%%
%%   ISO-code macros  \cmn{}, \jpn{} ... take a ROMANISATION and set it in
%%                    teal italic, optionally with a gloss: \cmn[dog]{gou3}.
%%   \zhs{} and friends take NATIVE SCRIPT and set it in teal upright in the
%%                    right regional face.  Script is never italicised.

\usepackage[c,e,f,j,k]{mtg2e}
\renewcommand{\mtcitestyle}[1]{\textcolor{teal}{\textit{#1}}}

% romanisations (mtg2e already gives \chn, \eng, \fra, \jpn, \kjpn and \kor)
\newcommand{\cmn}{\mtciteform}
\newcommand{\zh}{\mtciteform}
\newcommand{\yue}{\mtciteform}
\newcommand{\ko}{\mtciteform}
\newcommand{\vie}{\mtciteform}
\newcommand{\vi}{\mtciteform}
\newcommand{\tha}{\mtciteform}
\newcommand{\thaig}{\mtciteform}
\newcommand{\lao}{\mtciteform}
\newcommand{\ind}{\mtciteform}
\newcommand{\msa}{\mtciteform}
\newcommand{\zsm}{\mtciteform}
\newcommand{\ar}{\mtciteform}
\newcommand{\lex}[1]{\textbf{\mtcitestyle{#1}}}

% native script
\newcommand{\zhs}[1]{\mtplain{\textchinese{#1}}}
\newcommand{\zht}[1]{\mtplain{\textyue{#1}}}
\newcommand{\jps}[1]{\mtplain{\textjapanese{#1}}}
\newcommand{\kos}[1]{\mtplain{\textkorean{#1}}}
\newcommand{\ths}[1]{\mtplain{\textthai{#1}}}

%% ---- structure -----------------------------------------------------------
%% Sections are signposting, so give each one a page of its own showing where
%% we have got to, with the rest of the deck dimmed.  \overview puts the same
%% list up front; call it straight after \frame{\titlepage}.
\AtBeginSection[]{%
  \begin{frame}{Outline}
    \tableofcontents[currentsection]
  \end{frame}}

\newcommand{\overview}{%
  \begin{frame}{Overview}
    \tableofcontents
  \end{frame}}

%% ---- references ----------------------------------------------------------
\usepackage[round]{natbib}

%% ---- defaults ------------------------------------------------------------
%% The year lives here so the rollover is a one-line change.
\author[Francis Bond]{Francis Bond}
\date{2026}

%%   * \title, and \author where it is not just Francis Bond
%%   * anything a deck deliberately overrides afterwards -- sounds.tex, for
%%     instance, sets Andika as the body font because it is a phonetics deck
%%
%% Everything here can be overridden by the deck: %% shared.tex -- common preamble for the Linguistics of Asian Languages decks.
%%
%% \input this immediately after \documentclass{beamer}.  It fixes the theme,
%% the fonts for every script the course touches, the colour conventions and
%% the language shorthands in ONE place, so that a font fix made here reaches
%% every deck.
%%
%% Deliberately NOT here, because the decks genuinely differ:
%%   * the gb4e flavour -- gb4e, mygb4e and langsci-gb4e are not
%%     interchangeable, so each deck loads the one its examples are written
%%     for, straight after %% shared.tex -- common preamble for the Linguistics of Asian Languages decks.
%%
%% \input this immediately after \documentclass{beamer}.  It fixes the theme,
%% the fonts for every script the course touches, the colour conventions and
%% the language shorthands in ONE place, so that a font fix made here reaches
%% every deck.
%%
%% Deliberately NOT here, because the decks genuinely differ:
%%   * the gb4e flavour -- gb4e, mygb4e and langsci-gb4e are not
%%     interchangeable, so each deck loads the one its examples are written
%%     for, straight after %% shared.tex -- common preamble for the Linguistics of Asian Languages decks.
%%
%% \input this immediately after \documentclass{beamer}.  It fixes the theme,
%% the fonts for every script the course touches, the colour conventions and
%% the language shorthands in ONE place, so that a font fix made here reaches
%% every deck.
%%
%% Deliberately NOT here, because the decks genuinely differ:
%%   * the gb4e flavour -- gb4e, mygb4e and langsci-gb4e are not
%%     interchangeable, so each deck loads the one its examples are written
%%     for, straight after \input{shared}
%%   * \title, and \author where it is not just Francis Bond
%%   * anything a deck deliberately overrides afterwards -- sounds.tex, for
%%     instance, sets Andika as the body font because it is a phonetics deck
%%
%% Everything here can be overridden by the deck: \input{shared} first, then
%% deviate, and put a comment saying why.

%% ---- fonts -------------------------------------------------------------
%%
%% Two mechanisms, deliberately, because the decks use both:
%%
%%   xeCJK      picks the CJK font automatically from the character, so Han,
%%              kana and Hangul just work in running text.
%%   polyglossia gives \textchinese{}, \textjapanese{}, \textkorean{},
%%              \textyue{}, \textthai{} and \textarabic{} for when a specific
%%              regional face is wanted -- the Japanese and Traditional faces
%%              draw many shared characters differently from the Simplified
%%              one, and Arabic needs RTL.
%%
%% Han has no italic.  Map the italic shape onto the upright one for every
%% CJK face, or XeTeX fake-slants it whenever CJK lands inside \textit --
%% which it constantly does, because \mtcitestyle italicises.

\usepackage{fontspec}
\usepackage{xeCJK}
\usepackage{graphicx}
\usepackage{tikz}

% Body font.  Latin Modern Sans, beamer's default, is the wrong font for this
% course: it has no small caps (so \textsc{sg} in a gloss came out as
% lowercase "sg") and it is missing 33 of the characters these decks use --
% most of the IPA, the modifier letters and the super/subscripts -- which
% XeTeX then drops silently.
%
% Libertinus is the maintained successor to Linux Libertine/Biolinum.  Sans
% becomes the body font (beamer sets \sffamily), Serif is there for gloss
% lines, and it brings a real maths font with it.  It lives in TeX Live, so
% kpathsea finds it on any machine without a fontconfig install.
% Set the faces directly rather than via libertinus-otf, which also wants a
% mono font (AnonymousPro) that is not installed here.
% Ligatures={TeX,NoCommon}: keep ``---'' and friends, but switch OFF the fi/fl
% output ligatures, so that text copied out of the PDF comes back as plain
% ASCII.  Those ligatures are how the deck sources picked up "Pacific" and
% "geneDcally" in the first place.
\usefonttheme{professionalfonts}
\setmainfont{LibertinusSerif}[
  Extension		= .otf,
  UprightFont		= *-Regular,
  ItalicFont		= *-Italic,
  BoldFont		= *-Bold,
  BoldItalicFont	= *-BoldItalic,
  Ligatures		= {TeX,NoCommon}]
% Scale 0.95: Libertinus has a much larger x-height than Latin Modern, so at
% the nominal size it overflows frames that used to fit.  0.95 restores every
% deck to the page count it had before the switch.
\setsansfont{LibertinusSans}[
  Scale			= 0.95,
  Extension		= .otf,
  UprightFont		= *-Regular,
  ItalicFont		= *-Italic,
  BoldFont		= *-Bold,
  % Libertinus Sans ships no bold italic, so slant the bold rather than let
  % LaTeX silently substitute an upright bold.
  BoldItalicFont	= *-Bold,
  BoldItalicFeatures	= {FakeSlant = 0.2},
  Ligatures		= {TeX,NoCommon}]
\setmonofont{LibertinusMono-Regular}[Extension = .otf]

% XeTeX always loads face 0 of a TrueType Collection, and
% NotoSansCJK-Regular.ttc packs JP=0, KR=1, SC=2, TC=3, HK=4.  So asking for
% "Noto Sans CJK SC" -- as every deck did -- silently handed back the
% JAPANESE face, and all the Simplified Chinese in this course was being set
% with Japanese glyph shapes (骨, 直, 关 and many more differ).  FontIndex
% picks the right face, and the bare filename resolves without a path.
%
% Han has no italic.  Point ItalicFont at the same face so that CJK landing
% inside \textit -- which happens constantly, because \mtcitestyle
% italicises -- stays upright instead of drawing an "undefined shape".
\newcommand{\CJKface}[2]{%
  \newCJKfontfamily#1{NotoSansCJK-Regular.ttc}[%
    FontIndex = #2, ItalicFont = NotoSansCJK-Regular.ttc,
    ItalicFeatures = {FontIndex = #2}]}

\IfFontExistsTF{NotoSansCJK-Regular.ttc}{%
  \setCJKmainfont{NotoSansCJK-Regular.ttc}[%
    FontIndex = 2, ItalicFont = NotoSansCJK-Regular.ttc,
    ItalicFeatures = {FontIndex = 2}]
  \CJKface\chinesefontsf{2}	% Simplified
  \CJKface\yuefontsf{3}	% Traditional, used for Cantonese
  \CJKface\japanesefontsf{0}
  \CJKface\koreanfontsf{1}
}{%
  % No collection: fall back to the family names.  Everything still builds,
  % but the regional faces will all be whichever one fontconfig hands over.
  \setCJKmainfont{Noto Sans CJK SC}[ItalicFont = Noto Sans CJK SC]
  \newCJKfontfamily\chinesefontsf{Noto Sans CJK SC}[ItalicFont = Noto Sans CJK SC]
  \newCJKfontfamily\yuefontsf{Noto Sans CJK TC}[ItalicFont = Noto Sans CJK TC]
  \newCJKfontfamily\japanesefontsf{Noto Sans CJK JP}[ItalicFont = Noto Sans CJK JP]
  \newCJKfontfamily\koreanfontsf{Noto Sans CJK KR}[ItalicFont = Noto Sans CJK KR]
}
\setCJKsansfont{NotoSansCJK-Regular.ttc}[%
  FontIndex = 2, ItalicFont = NotoSansCJK-Regular.ttc,
  ItalicFeatures = {FontIndex = 2}]
% The mono faces live in the same collection: Mono JP=5, KR=6, SC=7, TC=8.
\setCJKmonofont{NotoSansCJK-Regular.ttc}[FontIndex = 7]

% The bopomofo tone marks sit in Spacing Modifier Letters, not the Bopomofo
% block, so xeCJK would leave them to Latin Modern -- which has no ˊ or ˋ.
\xeCJKDeclareCharClass{CJK}{"02C7, "02C9, "02CA, "02CB, "02D9}

% The reverse problem: xeCJK counts the curly quotes as CJK punctuation, so
% ' ' " " came out full-width in the CJK font, with CJK spacing around them
% and a line-break opportunity inside words like "doesn't".  In this course
% they are always English punctuation.
\xeCJKDeclareCharClass{Default}{"2018, "2019, "201C, "201D}

% IPA needs no font of its own any more: Biolinum covers IPA Extensions,
% Phonetic Extensions and the Spacing Modifier Letters.  The decks had grown
% four different answers to this (Andika, Noto Serif, Libertine, and nothing
% at all), so \ipa, \IPA and \silipa are all kept as names but now resolve
% to the body font.
\newfontfamily\ipafont{LibertinusSans}[
  Extension = .otf, UprightFont = *-Regular, ItalicFont = *-Italic,
  BoldFont = *-Bold, Ligatures = {TeX,NoCommon}]
\newcommand{\ipa}[1]{{\ipafont #1}}
\newcommand{\IPA}[1]{{\ipafont #1}}
\let\silipa\ipafont

% Used for the second line of gb4e glosses, where tone numbers and IPA turn
% up and Latin Modern runs out of superscripts.
\newfontfamily\Libertine{LibertinusSerif}[
  Extension = .otf, UprightFont = *-Regular, ItalicFont = *-Italic,
  BoldFont = *-Bold, BoldItalicFont = *-BoldItalic,
  Ligatures = {TeX,NoCommon}]

\usetheme{Madrid}
\usecolortheme{crane}

\usepackage{polyglossia}
\setdefaultlanguage{english}
\setotherlanguage{thai}
\newfontfamily\thaifont[Script=Thai]{Noto Sans Thai}
\newfontfamily\thaifontsf[Script=Thai]{Noto Sans Thai}
\setotherlanguage{arabic}
\newfontfamily\arabicfont[Script=Arabic]{Noto Naskh Arabic}
\newfontfamily\arabicfontsf[Script=Arabic]{Noto Naskh Arabic}

% polyglossia has no gloss file for these, so define the switches directly
% against the CJK families declared above.
\newcommand{\textchinese}[1]{{\chinesefontsf #1}}
\newcommand{\textjapanese}[1]{{\japanesefontsf #1}}
\newcommand{\textkorean}[1]{{\koreanfontsf #1}}
\newcommand{\textyue}[1]{{\yuefontsf #1}}

% Latin Modern has no mathematical angle brackets, which the writing systems
% deck uses for graphemes.
\usepackage{newunicodechar}
\newunicodechar{⟨}{‹}
\newunicodechar{⟩}{›}

%% ---- colour and emphasis ------------------------------------------------
\usepackage{xcolor}
\newcommand{\txx}[1]{\textcolor{blue}{\textbf{#1}}}
\let\emp\txx
\let\zz\txx
\newcommand{\mtplain}[1]{\textcolor{teal}{#1}}
\newcommand{\con}[1]{\textsc{#1}}

\usepackage[normalem]{ulem}
\newcommand{\ul}[1]{\uline{#1}}
\newcommand{\cll}[1]{\uline{#1}}

%% ---- linguistic examples ------------------------------------------------
%%
%% langsci-gb4e everywhere.  Plain gb4e makes _ and ^ active characters,
%% which breaks hyperref -- that is the one and only reason mygb4e existed, a
%% local copy of gb4e with exactly that hack commented out.  The Language
%% Science Press fork drops it, is maintained, and gives the \exfont /
%% \glossfont / \transfont hooks, so the local fork is no longer needed.
\usepackage{langsci-gb4e}

% Second gloss line -- transcription and IPA -- in the serif face.
% writing.tex and verbal-arts.tex override this afterwards.
\renewcommand{\eachwordtwo}{\Libertine}

%% ---- language shorthands ------------------------------------------------
%%
%% Two families, and the distinction matters:
%%
%%   ISO-code macros  \cmn{}, \jpn{} ... take a ROMANISATION and set it in
%%                    teal italic, optionally with a gloss: \cmn[dog]{gou3}.
%%   \zhs{} and friends take NATIVE SCRIPT and set it in teal upright in the
%%                    right regional face.  Script is never italicised.

\usepackage[c,e,f,j,k]{mtg2e}
\renewcommand{\mtcitestyle}[1]{\textcolor{teal}{\textit{#1}}}

% romanisations (mtg2e already gives \chn, \eng, \fra, \jpn, \kjpn and \kor)
\newcommand{\cmn}{\mtciteform}
\newcommand{\zh}{\mtciteform}
\newcommand{\yue}{\mtciteform}
\newcommand{\ko}{\mtciteform}
\newcommand{\vie}{\mtciteform}
\newcommand{\vi}{\mtciteform}
\newcommand{\tha}{\mtciteform}
\newcommand{\thaig}{\mtciteform}
\newcommand{\lao}{\mtciteform}
\newcommand{\ind}{\mtciteform}
\newcommand{\msa}{\mtciteform}
\newcommand{\zsm}{\mtciteform}
\newcommand{\ar}{\mtciteform}
\newcommand{\lex}[1]{\textbf{\mtcitestyle{#1}}}

% native script
\newcommand{\zhs}[1]{\mtplain{\textchinese{#1}}}
\newcommand{\zht}[1]{\mtplain{\textyue{#1}}}
\newcommand{\jps}[1]{\mtplain{\textjapanese{#1}}}
\newcommand{\kos}[1]{\mtplain{\textkorean{#1}}}
\newcommand{\ths}[1]{\mtplain{\textthai{#1}}}

%% ---- structure -----------------------------------------------------------
%% Sections are signposting, so give each one a page of its own showing where
%% we have got to, with the rest of the deck dimmed.  \overview puts the same
%% list up front; call it straight after \frame{\titlepage}.
\AtBeginSection[]{%
  \begin{frame}{Outline}
    \tableofcontents[currentsection]
  \end{frame}}

\newcommand{\overview}{%
  \begin{frame}{Overview}
    \tableofcontents
  \end{frame}}

%% ---- references ----------------------------------------------------------
\usepackage[round]{natbib}

%% ---- defaults ------------------------------------------------------------
%% The year lives here so the rollover is a one-line change.
\author[Francis Bond]{Francis Bond}
\date{2026}

%%   * \title, and \author where it is not just Francis Bond
%%   * anything a deck deliberately overrides afterwards -- sounds.tex, for
%%     instance, sets Andika as the body font because it is a phonetics deck
%%
%% Everything here can be overridden by the deck: %% shared.tex -- common preamble for the Linguistics of Asian Languages decks.
%%
%% \input this immediately after \documentclass{beamer}.  It fixes the theme,
%% the fonts for every script the course touches, the colour conventions and
%% the language shorthands in ONE place, so that a font fix made here reaches
%% every deck.
%%
%% Deliberately NOT here, because the decks genuinely differ:
%%   * the gb4e flavour -- gb4e, mygb4e and langsci-gb4e are not
%%     interchangeable, so each deck loads the one its examples are written
%%     for, straight after \input{shared}
%%   * \title, and \author where it is not just Francis Bond
%%   * anything a deck deliberately overrides afterwards -- sounds.tex, for
%%     instance, sets Andika as the body font because it is a phonetics deck
%%
%% Everything here can be overridden by the deck: \input{shared} first, then
%% deviate, and put a comment saying why.

%% ---- fonts -------------------------------------------------------------
%%
%% Two mechanisms, deliberately, because the decks use both:
%%
%%   xeCJK      picks the CJK font automatically from the character, so Han,
%%              kana and Hangul just work in running text.
%%   polyglossia gives \textchinese{}, \textjapanese{}, \textkorean{},
%%              \textyue{}, \textthai{} and \textarabic{} for when a specific
%%              regional face is wanted -- the Japanese and Traditional faces
%%              draw many shared characters differently from the Simplified
%%              one, and Arabic needs RTL.
%%
%% Han has no italic.  Map the italic shape onto the upright one for every
%% CJK face, or XeTeX fake-slants it whenever CJK lands inside \textit --
%% which it constantly does, because \mtcitestyle italicises.

\usepackage{fontspec}
\usepackage{xeCJK}
\usepackage{graphicx}
\usepackage{tikz}

% Body font.  Latin Modern Sans, beamer's default, is the wrong font for this
% course: it has no small caps (so \textsc{sg} in a gloss came out as
% lowercase "sg") and it is missing 33 of the characters these decks use --
% most of the IPA, the modifier letters and the super/subscripts -- which
% XeTeX then drops silently.
%
% Libertinus is the maintained successor to Linux Libertine/Biolinum.  Sans
% becomes the body font (beamer sets \sffamily), Serif is there for gloss
% lines, and it brings a real maths font with it.  It lives in TeX Live, so
% kpathsea finds it on any machine without a fontconfig install.
% Set the faces directly rather than via libertinus-otf, which also wants a
% mono font (AnonymousPro) that is not installed here.
% Ligatures={TeX,NoCommon}: keep ``---'' and friends, but switch OFF the fi/fl
% output ligatures, so that text copied out of the PDF comes back as plain
% ASCII.  Those ligatures are how the deck sources picked up "Pacific" and
% "geneDcally" in the first place.
\usefonttheme{professionalfonts}
\setmainfont{LibertinusSerif}[
  Extension		= .otf,
  UprightFont		= *-Regular,
  ItalicFont		= *-Italic,
  BoldFont		= *-Bold,
  BoldItalicFont	= *-BoldItalic,
  Ligatures		= {TeX,NoCommon}]
% Scale 0.95: Libertinus has a much larger x-height than Latin Modern, so at
% the nominal size it overflows frames that used to fit.  0.95 restores every
% deck to the page count it had before the switch.
\setsansfont{LibertinusSans}[
  Scale			= 0.95,
  Extension		= .otf,
  UprightFont		= *-Regular,
  ItalicFont		= *-Italic,
  BoldFont		= *-Bold,
  % Libertinus Sans ships no bold italic, so slant the bold rather than let
  % LaTeX silently substitute an upright bold.
  BoldItalicFont	= *-Bold,
  BoldItalicFeatures	= {FakeSlant = 0.2},
  Ligatures		= {TeX,NoCommon}]
\setmonofont{LibertinusMono-Regular}[Extension = .otf]

% XeTeX always loads face 0 of a TrueType Collection, and
% NotoSansCJK-Regular.ttc packs JP=0, KR=1, SC=2, TC=3, HK=4.  So asking for
% "Noto Sans CJK SC" -- as every deck did -- silently handed back the
% JAPANESE face, and all the Simplified Chinese in this course was being set
% with Japanese glyph shapes (骨, 直, 关 and many more differ).  FontIndex
% picks the right face, and the bare filename resolves without a path.
%
% Han has no italic.  Point ItalicFont at the same face so that CJK landing
% inside \textit -- which happens constantly, because \mtcitestyle
% italicises -- stays upright instead of drawing an "undefined shape".
\newcommand{\CJKface}[2]{%
  \newCJKfontfamily#1{NotoSansCJK-Regular.ttc}[%
    FontIndex = #2, ItalicFont = NotoSansCJK-Regular.ttc,
    ItalicFeatures = {FontIndex = #2}]}

\IfFontExistsTF{NotoSansCJK-Regular.ttc}{%
  \setCJKmainfont{NotoSansCJK-Regular.ttc}[%
    FontIndex = 2, ItalicFont = NotoSansCJK-Regular.ttc,
    ItalicFeatures = {FontIndex = 2}]
  \CJKface\chinesefontsf{2}	% Simplified
  \CJKface\yuefontsf{3}	% Traditional, used for Cantonese
  \CJKface\japanesefontsf{0}
  \CJKface\koreanfontsf{1}
}{%
  % No collection: fall back to the family names.  Everything still builds,
  % but the regional faces will all be whichever one fontconfig hands over.
  \setCJKmainfont{Noto Sans CJK SC}[ItalicFont = Noto Sans CJK SC]
  \newCJKfontfamily\chinesefontsf{Noto Sans CJK SC}[ItalicFont = Noto Sans CJK SC]
  \newCJKfontfamily\yuefontsf{Noto Sans CJK TC}[ItalicFont = Noto Sans CJK TC]
  \newCJKfontfamily\japanesefontsf{Noto Sans CJK JP}[ItalicFont = Noto Sans CJK JP]
  \newCJKfontfamily\koreanfontsf{Noto Sans CJK KR}[ItalicFont = Noto Sans CJK KR]
}
\setCJKsansfont{NotoSansCJK-Regular.ttc}[%
  FontIndex = 2, ItalicFont = NotoSansCJK-Regular.ttc,
  ItalicFeatures = {FontIndex = 2}]
% The mono faces live in the same collection: Mono JP=5, KR=6, SC=7, TC=8.
\setCJKmonofont{NotoSansCJK-Regular.ttc}[FontIndex = 7]

% The bopomofo tone marks sit in Spacing Modifier Letters, not the Bopomofo
% block, so xeCJK would leave them to Latin Modern -- which has no ˊ or ˋ.
\xeCJKDeclareCharClass{CJK}{"02C7, "02C9, "02CA, "02CB, "02D9}

% The reverse problem: xeCJK counts the curly quotes as CJK punctuation, so
% ' ' " " came out full-width in the CJK font, with CJK spacing around them
% and a line-break opportunity inside words like "doesn't".  In this course
% they are always English punctuation.
\xeCJKDeclareCharClass{Default}{"2018, "2019, "201C, "201D}

% IPA needs no font of its own any more: Biolinum covers IPA Extensions,
% Phonetic Extensions and the Spacing Modifier Letters.  The decks had grown
% four different answers to this (Andika, Noto Serif, Libertine, and nothing
% at all), so \ipa, \IPA and \silipa are all kept as names but now resolve
% to the body font.
\newfontfamily\ipafont{LibertinusSans}[
  Extension = .otf, UprightFont = *-Regular, ItalicFont = *-Italic,
  BoldFont = *-Bold, Ligatures = {TeX,NoCommon}]
\newcommand{\ipa}[1]{{\ipafont #1}}
\newcommand{\IPA}[1]{{\ipafont #1}}
\let\silipa\ipafont

% Used for the second line of gb4e glosses, where tone numbers and IPA turn
% up and Latin Modern runs out of superscripts.
\newfontfamily\Libertine{LibertinusSerif}[
  Extension = .otf, UprightFont = *-Regular, ItalicFont = *-Italic,
  BoldFont = *-Bold, BoldItalicFont = *-BoldItalic,
  Ligatures = {TeX,NoCommon}]

\usetheme{Madrid}
\usecolortheme{crane}

\usepackage{polyglossia}
\setdefaultlanguage{english}
\setotherlanguage{thai}
\newfontfamily\thaifont[Script=Thai]{Noto Sans Thai}
\newfontfamily\thaifontsf[Script=Thai]{Noto Sans Thai}
\setotherlanguage{arabic}
\newfontfamily\arabicfont[Script=Arabic]{Noto Naskh Arabic}
\newfontfamily\arabicfontsf[Script=Arabic]{Noto Naskh Arabic}

% polyglossia has no gloss file for these, so define the switches directly
% against the CJK families declared above.
\newcommand{\textchinese}[1]{{\chinesefontsf #1}}
\newcommand{\textjapanese}[1]{{\japanesefontsf #1}}
\newcommand{\textkorean}[1]{{\koreanfontsf #1}}
\newcommand{\textyue}[1]{{\yuefontsf #1}}

% Latin Modern has no mathematical angle brackets, which the writing systems
% deck uses for graphemes.
\usepackage{newunicodechar}
\newunicodechar{⟨}{‹}
\newunicodechar{⟩}{›}

%% ---- colour and emphasis ------------------------------------------------
\usepackage{xcolor}
\newcommand{\txx}[1]{\textcolor{blue}{\textbf{#1}}}
\let\emp\txx
\let\zz\txx
\newcommand{\mtplain}[1]{\textcolor{teal}{#1}}
\newcommand{\con}[1]{\textsc{#1}}

\usepackage[normalem]{ulem}
\newcommand{\ul}[1]{\uline{#1}}
\newcommand{\cll}[1]{\uline{#1}}

%% ---- linguistic examples ------------------------------------------------
%%
%% langsci-gb4e everywhere.  Plain gb4e makes _ and ^ active characters,
%% which breaks hyperref -- that is the one and only reason mygb4e existed, a
%% local copy of gb4e with exactly that hack commented out.  The Language
%% Science Press fork drops it, is maintained, and gives the \exfont /
%% \glossfont / \transfont hooks, so the local fork is no longer needed.
\usepackage{langsci-gb4e}

% Second gloss line -- transcription and IPA -- in the serif face.
% writing.tex and verbal-arts.tex override this afterwards.
\renewcommand{\eachwordtwo}{\Libertine}

%% ---- language shorthands ------------------------------------------------
%%
%% Two families, and the distinction matters:
%%
%%   ISO-code macros  \cmn{}, \jpn{} ... take a ROMANISATION and set it in
%%                    teal italic, optionally with a gloss: \cmn[dog]{gou3}.
%%   \zhs{} and friends take NATIVE SCRIPT and set it in teal upright in the
%%                    right regional face.  Script is never italicised.

\usepackage[c,e,f,j,k]{mtg2e}
\renewcommand{\mtcitestyle}[1]{\textcolor{teal}{\textit{#1}}}

% romanisations (mtg2e already gives \chn, \eng, \fra, \jpn, \kjpn and \kor)
\newcommand{\cmn}{\mtciteform}
\newcommand{\zh}{\mtciteform}
\newcommand{\yue}{\mtciteform}
\newcommand{\ko}{\mtciteform}
\newcommand{\vie}{\mtciteform}
\newcommand{\vi}{\mtciteform}
\newcommand{\tha}{\mtciteform}
\newcommand{\thaig}{\mtciteform}
\newcommand{\lao}{\mtciteform}
\newcommand{\ind}{\mtciteform}
\newcommand{\msa}{\mtciteform}
\newcommand{\zsm}{\mtciteform}
\newcommand{\ar}{\mtciteform}
\newcommand{\lex}[1]{\textbf{\mtcitestyle{#1}}}

% native script
\newcommand{\zhs}[1]{\mtplain{\textchinese{#1}}}
\newcommand{\zht}[1]{\mtplain{\textyue{#1}}}
\newcommand{\jps}[1]{\mtplain{\textjapanese{#1}}}
\newcommand{\kos}[1]{\mtplain{\textkorean{#1}}}
\newcommand{\ths}[1]{\mtplain{\textthai{#1}}}

%% ---- structure -----------------------------------------------------------
%% Sections are signposting, so give each one a page of its own showing where
%% we have got to, with the rest of the deck dimmed.  \overview puts the same
%% list up front; call it straight after \frame{\titlepage}.
\AtBeginSection[]{%
  \begin{frame}{Outline}
    \tableofcontents[currentsection]
  \end{frame}}

\newcommand{\overview}{%
  \begin{frame}{Overview}
    \tableofcontents
  \end{frame}}

%% ---- references ----------------------------------------------------------
\usepackage[round]{natbib}

%% ---- defaults ------------------------------------------------------------
%% The year lives here so the rollover is a one-line change.
\author[Francis Bond]{Francis Bond}
\date{2026}
 first, then
%% deviate, and put a comment saying why.

%% ---- fonts -------------------------------------------------------------
%%
%% Two mechanisms, deliberately, because the decks use both:
%%
%%   xeCJK      picks the CJK font automatically from the character, so Han,
%%              kana and Hangul just work in running text.
%%   polyglossia gives \textchinese{}, \textjapanese{}, \textkorean{},
%%              \textyue{}, \textthai{} and \textarabic{} for when a specific
%%              regional face is wanted -- the Japanese and Traditional faces
%%              draw many shared characters differently from the Simplified
%%              one, and Arabic needs RTL.
%%
%% Han has no italic.  Map the italic shape onto the upright one for every
%% CJK face, or XeTeX fake-slants it whenever CJK lands inside \textit --
%% which it constantly does, because \mtcitestyle italicises.

\usepackage{fontspec}
\usepackage{xeCJK}
\usepackage{graphicx}
\usepackage{tikz}

% Body font.  Latin Modern Sans, beamer's default, is the wrong font for this
% course: it has no small caps (so \textsc{sg} in a gloss came out as
% lowercase "sg") and it is missing 33 of the characters these decks use --
% most of the IPA, the modifier letters and the super/subscripts -- which
% XeTeX then drops silently.
%
% Libertinus is the maintained successor to Linux Libertine/Biolinum.  Sans
% becomes the body font (beamer sets \sffamily), Serif is there for gloss
% lines, and it brings a real maths font with it.  It lives in TeX Live, so
% kpathsea finds it on any machine without a fontconfig install.
% Set the faces directly rather than via libertinus-otf, which also wants a
% mono font (AnonymousPro) that is not installed here.
% Ligatures={TeX,NoCommon}: keep ``---'' and friends, but switch OFF the fi/fl
% output ligatures, so that text copied out of the PDF comes back as plain
% ASCII.  Those ligatures are how the deck sources picked up "Pacific" and
% "geneDcally" in the first place.
\usefonttheme{professionalfonts}
\setmainfont{LibertinusSerif}[
  Extension		= .otf,
  UprightFont		= *-Regular,
  ItalicFont		= *-Italic,
  BoldFont		= *-Bold,
  BoldItalicFont	= *-BoldItalic,
  Ligatures		= {TeX,NoCommon}]
% Scale 0.95: Libertinus has a much larger x-height than Latin Modern, so at
% the nominal size it overflows frames that used to fit.  0.95 restores every
% deck to the page count it had before the switch.
\setsansfont{LibertinusSans}[
  Scale			= 0.95,
  Extension		= .otf,
  UprightFont		= *-Regular,
  ItalicFont		= *-Italic,
  BoldFont		= *-Bold,
  % Libertinus Sans ships no bold italic, so slant the bold rather than let
  % LaTeX silently substitute an upright bold.
  BoldItalicFont	= *-Bold,
  BoldItalicFeatures	= {FakeSlant = 0.2},
  Ligatures		= {TeX,NoCommon}]
\setmonofont{LibertinusMono-Regular}[Extension = .otf]

% XeTeX always loads face 0 of a TrueType Collection, and
% NotoSansCJK-Regular.ttc packs JP=0, KR=1, SC=2, TC=3, HK=4.  So asking for
% "Noto Sans CJK SC" -- as every deck did -- silently handed back the
% JAPANESE face, and all the Simplified Chinese in this course was being set
% with Japanese glyph shapes (骨, 直, 关 and many more differ).  FontIndex
% picks the right face, and the bare filename resolves without a path.
%
% Han has no italic.  Point ItalicFont at the same face so that CJK landing
% inside \textit -- which happens constantly, because \mtcitestyle
% italicises -- stays upright instead of drawing an "undefined shape".
\newcommand{\CJKface}[2]{%
  \newCJKfontfamily#1{NotoSansCJK-Regular.ttc}[%
    FontIndex = #2, ItalicFont = NotoSansCJK-Regular.ttc,
    ItalicFeatures = {FontIndex = #2}]}

\IfFontExistsTF{NotoSansCJK-Regular.ttc}{%
  \setCJKmainfont{NotoSansCJK-Regular.ttc}[%
    FontIndex = 2, ItalicFont = NotoSansCJK-Regular.ttc,
    ItalicFeatures = {FontIndex = 2}]
  \CJKface\chinesefontsf{2}	% Simplified
  \CJKface\yuefontsf{3}	% Traditional, used for Cantonese
  \CJKface\japanesefontsf{0}
  \CJKface\koreanfontsf{1}
}{%
  % No collection: fall back to the family names.  Everything still builds,
  % but the regional faces will all be whichever one fontconfig hands over.
  \setCJKmainfont{Noto Sans CJK SC}[ItalicFont = Noto Sans CJK SC]
  \newCJKfontfamily\chinesefontsf{Noto Sans CJK SC}[ItalicFont = Noto Sans CJK SC]
  \newCJKfontfamily\yuefontsf{Noto Sans CJK TC}[ItalicFont = Noto Sans CJK TC]
  \newCJKfontfamily\japanesefontsf{Noto Sans CJK JP}[ItalicFont = Noto Sans CJK JP]
  \newCJKfontfamily\koreanfontsf{Noto Sans CJK KR}[ItalicFont = Noto Sans CJK KR]
}
\setCJKsansfont{NotoSansCJK-Regular.ttc}[%
  FontIndex = 2, ItalicFont = NotoSansCJK-Regular.ttc,
  ItalicFeatures = {FontIndex = 2}]
% The mono faces live in the same collection: Mono JP=5, KR=6, SC=7, TC=8.
\setCJKmonofont{NotoSansCJK-Regular.ttc}[FontIndex = 7]

% The bopomofo tone marks sit in Spacing Modifier Letters, not the Bopomofo
% block, so xeCJK would leave them to Latin Modern -- which has no ˊ or ˋ.
\xeCJKDeclareCharClass{CJK}{"02C7, "02C9, "02CA, "02CB, "02D9}

% The reverse problem: xeCJK counts the curly quotes as CJK punctuation, so
% ' ' " " came out full-width in the CJK font, with CJK spacing around them
% and a line-break opportunity inside words like "doesn't".  In this course
% they are always English punctuation.
\xeCJKDeclareCharClass{Default}{"2018, "2019, "201C, "201D}

% IPA needs no font of its own any more: Biolinum covers IPA Extensions,
% Phonetic Extensions and the Spacing Modifier Letters.  The decks had grown
% four different answers to this (Andika, Noto Serif, Libertine, and nothing
% at all), so \ipa, \IPA and \silipa are all kept as names but now resolve
% to the body font.
\newfontfamily\ipafont{LibertinusSans}[
  Extension = .otf, UprightFont = *-Regular, ItalicFont = *-Italic,
  BoldFont = *-Bold, Ligatures = {TeX,NoCommon}]
\newcommand{\ipa}[1]{{\ipafont #1}}
\newcommand{\IPA}[1]{{\ipafont #1}}
\let\silipa\ipafont

% Used for the second line of gb4e glosses, where tone numbers and IPA turn
% up and Latin Modern runs out of superscripts.
\newfontfamily\Libertine{LibertinusSerif}[
  Extension = .otf, UprightFont = *-Regular, ItalicFont = *-Italic,
  BoldFont = *-Bold, BoldItalicFont = *-BoldItalic,
  Ligatures = {TeX,NoCommon}]

\usetheme{Madrid}
\usecolortheme{crane}

\usepackage{polyglossia}
\setdefaultlanguage{english}
\setotherlanguage{thai}
\newfontfamily\thaifont[Script=Thai]{Noto Sans Thai}
\newfontfamily\thaifontsf[Script=Thai]{Noto Sans Thai}
\setotherlanguage{arabic}
\newfontfamily\arabicfont[Script=Arabic]{Noto Naskh Arabic}
\newfontfamily\arabicfontsf[Script=Arabic]{Noto Naskh Arabic}

% polyglossia has no gloss file for these, so define the switches directly
% against the CJK families declared above.
\newcommand{\textchinese}[1]{{\chinesefontsf #1}}
\newcommand{\textjapanese}[1]{{\japanesefontsf #1}}
\newcommand{\textkorean}[1]{{\koreanfontsf #1}}
\newcommand{\textyue}[1]{{\yuefontsf #1}}

% Latin Modern has no mathematical angle brackets, which the writing systems
% deck uses for graphemes.
\usepackage{newunicodechar}
\newunicodechar{⟨}{‹}
\newunicodechar{⟩}{›}

%% ---- colour and emphasis ------------------------------------------------
\usepackage{xcolor}
\newcommand{\txx}[1]{\textcolor{blue}{\textbf{#1}}}
\let\emp\txx
\let\zz\txx
\newcommand{\mtplain}[1]{\textcolor{teal}{#1}}
\newcommand{\con}[1]{\textsc{#1}}

\usepackage[normalem]{ulem}
\newcommand{\ul}[1]{\uline{#1}}
\newcommand{\cll}[1]{\uline{#1}}

%% ---- linguistic examples ------------------------------------------------
%%
%% langsci-gb4e everywhere.  Plain gb4e makes _ and ^ active characters,
%% which breaks hyperref -- that is the one and only reason mygb4e existed, a
%% local copy of gb4e with exactly that hack commented out.  The Language
%% Science Press fork drops it, is maintained, and gives the \exfont /
%% \glossfont / \transfont hooks, so the local fork is no longer needed.
\usepackage{langsci-gb4e}

% Second gloss line -- transcription and IPA -- in the serif face.
% writing.tex and verbal-arts.tex override this afterwards.
\renewcommand{\eachwordtwo}{\Libertine}

%% ---- language shorthands ------------------------------------------------
%%
%% Two families, and the distinction matters:
%%
%%   ISO-code macros  \cmn{}, \jpn{} ... take a ROMANISATION and set it in
%%                    teal italic, optionally with a gloss: \cmn[dog]{gou3}.
%%   \zhs{} and friends take NATIVE SCRIPT and set it in teal upright in the
%%                    right regional face.  Script is never italicised.

\usepackage[c,e,f,j,k]{mtg2e}
\renewcommand{\mtcitestyle}[1]{\textcolor{teal}{\textit{#1}}}

% romanisations (mtg2e already gives \chn, \eng, \fra, \jpn, \kjpn and \kor)
\newcommand{\cmn}{\mtciteform}
\newcommand{\zh}{\mtciteform}
\newcommand{\yue}{\mtciteform}
\newcommand{\ko}{\mtciteform}
\newcommand{\vie}{\mtciteform}
\newcommand{\vi}{\mtciteform}
\newcommand{\tha}{\mtciteform}
\newcommand{\thaig}{\mtciteform}
\newcommand{\lao}{\mtciteform}
\newcommand{\ind}{\mtciteform}
\newcommand{\msa}{\mtciteform}
\newcommand{\zsm}{\mtciteform}
\newcommand{\ar}{\mtciteform}
\newcommand{\lex}[1]{\textbf{\mtcitestyle{#1}}}

% native script
\newcommand{\zhs}[1]{\mtplain{\textchinese{#1}}}
\newcommand{\zht}[1]{\mtplain{\textyue{#1}}}
\newcommand{\jps}[1]{\mtplain{\textjapanese{#1}}}
\newcommand{\kos}[1]{\mtplain{\textkorean{#1}}}
\newcommand{\ths}[1]{\mtplain{\textthai{#1}}}

%% ---- structure -----------------------------------------------------------
%% Sections are signposting, so give each one a page of its own showing where
%% we have got to, with the rest of the deck dimmed.  \overview puts the same
%% list up front; call it straight after \frame{\titlepage}.
\AtBeginSection[]{%
  \begin{frame}{Outline}
    \tableofcontents[currentsection]
  \end{frame}}

\newcommand{\overview}{%
  \begin{frame}{Overview}
    \tableofcontents
  \end{frame}}

%% ---- references ----------------------------------------------------------
\usepackage[round]{natbib}

%% ---- defaults ------------------------------------------------------------
%% The year lives here so the rollover is a one-line change.
\author[Francis Bond]{Francis Bond}
\date{2026}

%%   * \title, and \author where it is not just Francis Bond
%%   * anything a deck deliberately overrides afterwards -- sounds.tex, for
%%     instance, sets Andika as the body font because it is a phonetics deck
%%
%% Everything here can be overridden by the deck: %% shared.tex -- common preamble for the Linguistics of Asian Languages decks.
%%
%% \input this immediately after \documentclass{beamer}.  It fixes the theme,
%% the fonts for every script the course touches, the colour conventions and
%% the language shorthands in ONE place, so that a font fix made here reaches
%% every deck.
%%
%% Deliberately NOT here, because the decks genuinely differ:
%%   * the gb4e flavour -- gb4e, mygb4e and langsci-gb4e are not
%%     interchangeable, so each deck loads the one its examples are written
%%     for, straight after %% shared.tex -- common preamble for the Linguistics of Asian Languages decks.
%%
%% \input this immediately after \documentclass{beamer}.  It fixes the theme,
%% the fonts for every script the course touches, the colour conventions and
%% the language shorthands in ONE place, so that a font fix made here reaches
%% every deck.
%%
%% Deliberately NOT here, because the decks genuinely differ:
%%   * the gb4e flavour -- gb4e, mygb4e and langsci-gb4e are not
%%     interchangeable, so each deck loads the one its examples are written
%%     for, straight after \input{shared}
%%   * \title, and \author where it is not just Francis Bond
%%   * anything a deck deliberately overrides afterwards -- sounds.tex, for
%%     instance, sets Andika as the body font because it is a phonetics deck
%%
%% Everything here can be overridden by the deck: \input{shared} first, then
%% deviate, and put a comment saying why.

%% ---- fonts -------------------------------------------------------------
%%
%% Two mechanisms, deliberately, because the decks use both:
%%
%%   xeCJK      picks the CJK font automatically from the character, so Han,
%%              kana and Hangul just work in running text.
%%   polyglossia gives \textchinese{}, \textjapanese{}, \textkorean{},
%%              \textyue{}, \textthai{} and \textarabic{} for when a specific
%%              regional face is wanted -- the Japanese and Traditional faces
%%              draw many shared characters differently from the Simplified
%%              one, and Arabic needs RTL.
%%
%% Han has no italic.  Map the italic shape onto the upright one for every
%% CJK face, or XeTeX fake-slants it whenever CJK lands inside \textit --
%% which it constantly does, because \mtcitestyle italicises.

\usepackage{fontspec}
\usepackage{xeCJK}
\usepackage{graphicx}
\usepackage{tikz}

% Body font.  Latin Modern Sans, beamer's default, is the wrong font for this
% course: it has no small caps (so \textsc{sg} in a gloss came out as
% lowercase "sg") and it is missing 33 of the characters these decks use --
% most of the IPA, the modifier letters and the super/subscripts -- which
% XeTeX then drops silently.
%
% Libertinus is the maintained successor to Linux Libertine/Biolinum.  Sans
% becomes the body font (beamer sets \sffamily), Serif is there for gloss
% lines, and it brings a real maths font with it.  It lives in TeX Live, so
% kpathsea finds it on any machine without a fontconfig install.
% Set the faces directly rather than via libertinus-otf, which also wants a
% mono font (AnonymousPro) that is not installed here.
% Ligatures={TeX,NoCommon}: keep ``---'' and friends, but switch OFF the fi/fl
% output ligatures, so that text copied out of the PDF comes back as plain
% ASCII.  Those ligatures are how the deck sources picked up "Pacific" and
% "geneDcally" in the first place.
\usefonttheme{professionalfonts}
\setmainfont{LibertinusSerif}[
  Extension		= .otf,
  UprightFont		= *-Regular,
  ItalicFont		= *-Italic,
  BoldFont		= *-Bold,
  BoldItalicFont	= *-BoldItalic,
  Ligatures		= {TeX,NoCommon}]
% Scale 0.95: Libertinus has a much larger x-height than Latin Modern, so at
% the nominal size it overflows frames that used to fit.  0.95 restores every
% deck to the page count it had before the switch.
\setsansfont{LibertinusSans}[
  Scale			= 0.95,
  Extension		= .otf,
  UprightFont		= *-Regular,
  ItalicFont		= *-Italic,
  BoldFont		= *-Bold,
  % Libertinus Sans ships no bold italic, so slant the bold rather than let
  % LaTeX silently substitute an upright bold.
  BoldItalicFont	= *-Bold,
  BoldItalicFeatures	= {FakeSlant = 0.2},
  Ligatures		= {TeX,NoCommon}]
\setmonofont{LibertinusMono-Regular}[Extension = .otf]

% XeTeX always loads face 0 of a TrueType Collection, and
% NotoSansCJK-Regular.ttc packs JP=0, KR=1, SC=2, TC=3, HK=4.  So asking for
% "Noto Sans CJK SC" -- as every deck did -- silently handed back the
% JAPANESE face, and all the Simplified Chinese in this course was being set
% with Japanese glyph shapes (骨, 直, 关 and many more differ).  FontIndex
% picks the right face, and the bare filename resolves without a path.
%
% Han has no italic.  Point ItalicFont at the same face so that CJK landing
% inside \textit -- which happens constantly, because \mtcitestyle
% italicises -- stays upright instead of drawing an "undefined shape".
\newcommand{\CJKface}[2]{%
  \newCJKfontfamily#1{NotoSansCJK-Regular.ttc}[%
    FontIndex = #2, ItalicFont = NotoSansCJK-Regular.ttc,
    ItalicFeatures = {FontIndex = #2}]}

\IfFontExistsTF{NotoSansCJK-Regular.ttc}{%
  \setCJKmainfont{NotoSansCJK-Regular.ttc}[%
    FontIndex = 2, ItalicFont = NotoSansCJK-Regular.ttc,
    ItalicFeatures = {FontIndex = 2}]
  \CJKface\chinesefontsf{2}	% Simplified
  \CJKface\yuefontsf{3}	% Traditional, used for Cantonese
  \CJKface\japanesefontsf{0}
  \CJKface\koreanfontsf{1}
}{%
  % No collection: fall back to the family names.  Everything still builds,
  % but the regional faces will all be whichever one fontconfig hands over.
  \setCJKmainfont{Noto Sans CJK SC}[ItalicFont = Noto Sans CJK SC]
  \newCJKfontfamily\chinesefontsf{Noto Sans CJK SC}[ItalicFont = Noto Sans CJK SC]
  \newCJKfontfamily\yuefontsf{Noto Sans CJK TC}[ItalicFont = Noto Sans CJK TC]
  \newCJKfontfamily\japanesefontsf{Noto Sans CJK JP}[ItalicFont = Noto Sans CJK JP]
  \newCJKfontfamily\koreanfontsf{Noto Sans CJK KR}[ItalicFont = Noto Sans CJK KR]
}
\setCJKsansfont{NotoSansCJK-Regular.ttc}[%
  FontIndex = 2, ItalicFont = NotoSansCJK-Regular.ttc,
  ItalicFeatures = {FontIndex = 2}]
% The mono faces live in the same collection: Mono JP=5, KR=6, SC=7, TC=8.
\setCJKmonofont{NotoSansCJK-Regular.ttc}[FontIndex = 7]

% The bopomofo tone marks sit in Spacing Modifier Letters, not the Bopomofo
% block, so xeCJK would leave them to Latin Modern -- which has no ˊ or ˋ.
\xeCJKDeclareCharClass{CJK}{"02C7, "02C9, "02CA, "02CB, "02D9}

% The reverse problem: xeCJK counts the curly quotes as CJK punctuation, so
% ' ' " " came out full-width in the CJK font, with CJK spacing around them
% and a line-break opportunity inside words like "doesn't".  In this course
% they are always English punctuation.
\xeCJKDeclareCharClass{Default}{"2018, "2019, "201C, "201D}

% IPA needs no font of its own any more: Biolinum covers IPA Extensions,
% Phonetic Extensions and the Spacing Modifier Letters.  The decks had grown
% four different answers to this (Andika, Noto Serif, Libertine, and nothing
% at all), so \ipa, \IPA and \silipa are all kept as names but now resolve
% to the body font.
\newfontfamily\ipafont{LibertinusSans}[
  Extension = .otf, UprightFont = *-Regular, ItalicFont = *-Italic,
  BoldFont = *-Bold, Ligatures = {TeX,NoCommon}]
\newcommand{\ipa}[1]{{\ipafont #1}}
\newcommand{\IPA}[1]{{\ipafont #1}}
\let\silipa\ipafont

% Used for the second line of gb4e glosses, where tone numbers and IPA turn
% up and Latin Modern runs out of superscripts.
\newfontfamily\Libertine{LibertinusSerif}[
  Extension = .otf, UprightFont = *-Regular, ItalicFont = *-Italic,
  BoldFont = *-Bold, BoldItalicFont = *-BoldItalic,
  Ligatures = {TeX,NoCommon}]

\usetheme{Madrid}
\usecolortheme{crane}

\usepackage{polyglossia}
\setdefaultlanguage{english}
\setotherlanguage{thai}
\newfontfamily\thaifont[Script=Thai]{Noto Sans Thai}
\newfontfamily\thaifontsf[Script=Thai]{Noto Sans Thai}
\setotherlanguage{arabic}
\newfontfamily\arabicfont[Script=Arabic]{Noto Naskh Arabic}
\newfontfamily\arabicfontsf[Script=Arabic]{Noto Naskh Arabic}

% polyglossia has no gloss file for these, so define the switches directly
% against the CJK families declared above.
\newcommand{\textchinese}[1]{{\chinesefontsf #1}}
\newcommand{\textjapanese}[1]{{\japanesefontsf #1}}
\newcommand{\textkorean}[1]{{\koreanfontsf #1}}
\newcommand{\textyue}[1]{{\yuefontsf #1}}

% Latin Modern has no mathematical angle brackets, which the writing systems
% deck uses for graphemes.
\usepackage{newunicodechar}
\newunicodechar{⟨}{‹}
\newunicodechar{⟩}{›}

%% ---- colour and emphasis ------------------------------------------------
\usepackage{xcolor}
\newcommand{\txx}[1]{\textcolor{blue}{\textbf{#1}}}
\let\emp\txx
\let\zz\txx
\newcommand{\mtplain}[1]{\textcolor{teal}{#1}}
\newcommand{\con}[1]{\textsc{#1}}

\usepackage[normalem]{ulem}
\newcommand{\ul}[1]{\uline{#1}}
\newcommand{\cll}[1]{\uline{#1}}

%% ---- linguistic examples ------------------------------------------------
%%
%% langsci-gb4e everywhere.  Plain gb4e makes _ and ^ active characters,
%% which breaks hyperref -- that is the one and only reason mygb4e existed, a
%% local copy of gb4e with exactly that hack commented out.  The Language
%% Science Press fork drops it, is maintained, and gives the \exfont /
%% \glossfont / \transfont hooks, so the local fork is no longer needed.
\usepackage{langsci-gb4e}

% Second gloss line -- transcription and IPA -- in the serif face.
% writing.tex and verbal-arts.tex override this afterwards.
\renewcommand{\eachwordtwo}{\Libertine}

%% ---- language shorthands ------------------------------------------------
%%
%% Two families, and the distinction matters:
%%
%%   ISO-code macros  \cmn{}, \jpn{} ... take a ROMANISATION and set it in
%%                    teal italic, optionally with a gloss: \cmn[dog]{gou3}.
%%   \zhs{} and friends take NATIVE SCRIPT and set it in teal upright in the
%%                    right regional face.  Script is never italicised.

\usepackage[c,e,f,j,k]{mtg2e}
\renewcommand{\mtcitestyle}[1]{\textcolor{teal}{\textit{#1}}}

% romanisations (mtg2e already gives \chn, \eng, \fra, \jpn, \kjpn and \kor)
\newcommand{\cmn}{\mtciteform}
\newcommand{\zh}{\mtciteform}
\newcommand{\yue}{\mtciteform}
\newcommand{\ko}{\mtciteform}
\newcommand{\vie}{\mtciteform}
\newcommand{\vi}{\mtciteform}
\newcommand{\tha}{\mtciteform}
\newcommand{\thaig}{\mtciteform}
\newcommand{\lao}{\mtciteform}
\newcommand{\ind}{\mtciteform}
\newcommand{\msa}{\mtciteform}
\newcommand{\zsm}{\mtciteform}
\newcommand{\ar}{\mtciteform}
\newcommand{\lex}[1]{\textbf{\mtcitestyle{#1}}}

% native script
\newcommand{\zhs}[1]{\mtplain{\textchinese{#1}}}
\newcommand{\zht}[1]{\mtplain{\textyue{#1}}}
\newcommand{\jps}[1]{\mtplain{\textjapanese{#1}}}
\newcommand{\kos}[1]{\mtplain{\textkorean{#1}}}
\newcommand{\ths}[1]{\mtplain{\textthai{#1}}}

%% ---- structure -----------------------------------------------------------
%% Sections are signposting, so give each one a page of its own showing where
%% we have got to, with the rest of the deck dimmed.  \overview puts the same
%% list up front; call it straight after \frame{\titlepage}.
\AtBeginSection[]{%
  \begin{frame}{Outline}
    \tableofcontents[currentsection]
  \end{frame}}

\newcommand{\overview}{%
  \begin{frame}{Overview}
    \tableofcontents
  \end{frame}}

%% ---- references ----------------------------------------------------------
\usepackage[round]{natbib}

%% ---- defaults ------------------------------------------------------------
%% The year lives here so the rollover is a one-line change.
\author[Francis Bond]{Francis Bond}
\date{2026}

%%   * \title, and \author where it is not just Francis Bond
%%   * anything a deck deliberately overrides afterwards -- sounds.tex, for
%%     instance, sets Andika as the body font because it is a phonetics deck
%%
%% Everything here can be overridden by the deck: %% shared.tex -- common preamble for the Linguistics of Asian Languages decks.
%%
%% \input this immediately after \documentclass{beamer}.  It fixes the theme,
%% the fonts for every script the course touches, the colour conventions and
%% the language shorthands in ONE place, so that a font fix made here reaches
%% every deck.
%%
%% Deliberately NOT here, because the decks genuinely differ:
%%   * the gb4e flavour -- gb4e, mygb4e and langsci-gb4e are not
%%     interchangeable, so each deck loads the one its examples are written
%%     for, straight after \input{shared}
%%   * \title, and \author where it is not just Francis Bond
%%   * anything a deck deliberately overrides afterwards -- sounds.tex, for
%%     instance, sets Andika as the body font because it is a phonetics deck
%%
%% Everything here can be overridden by the deck: \input{shared} first, then
%% deviate, and put a comment saying why.

%% ---- fonts -------------------------------------------------------------
%%
%% Two mechanisms, deliberately, because the decks use both:
%%
%%   xeCJK      picks the CJK font automatically from the character, so Han,
%%              kana and Hangul just work in running text.
%%   polyglossia gives \textchinese{}, \textjapanese{}, \textkorean{},
%%              \textyue{}, \textthai{} and \textarabic{} for when a specific
%%              regional face is wanted -- the Japanese and Traditional faces
%%              draw many shared characters differently from the Simplified
%%              one, and Arabic needs RTL.
%%
%% Han has no italic.  Map the italic shape onto the upright one for every
%% CJK face, or XeTeX fake-slants it whenever CJK lands inside \textit --
%% which it constantly does, because \mtcitestyle italicises.

\usepackage{fontspec}
\usepackage{xeCJK}
\usepackage{graphicx}
\usepackage{tikz}

% Body font.  Latin Modern Sans, beamer's default, is the wrong font for this
% course: it has no small caps (so \textsc{sg} in a gloss came out as
% lowercase "sg") and it is missing 33 of the characters these decks use --
% most of the IPA, the modifier letters and the super/subscripts -- which
% XeTeX then drops silently.
%
% Libertinus is the maintained successor to Linux Libertine/Biolinum.  Sans
% becomes the body font (beamer sets \sffamily), Serif is there for gloss
% lines, and it brings a real maths font with it.  It lives in TeX Live, so
% kpathsea finds it on any machine without a fontconfig install.
% Set the faces directly rather than via libertinus-otf, which also wants a
% mono font (AnonymousPro) that is not installed here.
% Ligatures={TeX,NoCommon}: keep ``---'' and friends, but switch OFF the fi/fl
% output ligatures, so that text copied out of the PDF comes back as plain
% ASCII.  Those ligatures are how the deck sources picked up "Pacific" and
% "geneDcally" in the first place.
\usefonttheme{professionalfonts}
\setmainfont{LibertinusSerif}[
  Extension		= .otf,
  UprightFont		= *-Regular,
  ItalicFont		= *-Italic,
  BoldFont		= *-Bold,
  BoldItalicFont	= *-BoldItalic,
  Ligatures		= {TeX,NoCommon}]
% Scale 0.95: Libertinus has a much larger x-height than Latin Modern, so at
% the nominal size it overflows frames that used to fit.  0.95 restores every
% deck to the page count it had before the switch.
\setsansfont{LibertinusSans}[
  Scale			= 0.95,
  Extension		= .otf,
  UprightFont		= *-Regular,
  ItalicFont		= *-Italic,
  BoldFont		= *-Bold,
  % Libertinus Sans ships no bold italic, so slant the bold rather than let
  % LaTeX silently substitute an upright bold.
  BoldItalicFont	= *-Bold,
  BoldItalicFeatures	= {FakeSlant = 0.2},
  Ligatures		= {TeX,NoCommon}]
\setmonofont{LibertinusMono-Regular}[Extension = .otf]

% XeTeX always loads face 0 of a TrueType Collection, and
% NotoSansCJK-Regular.ttc packs JP=0, KR=1, SC=2, TC=3, HK=4.  So asking for
% "Noto Sans CJK SC" -- as every deck did -- silently handed back the
% JAPANESE face, and all the Simplified Chinese in this course was being set
% with Japanese glyph shapes (骨, 直, 关 and many more differ).  FontIndex
% picks the right face, and the bare filename resolves without a path.
%
% Han has no italic.  Point ItalicFont at the same face so that CJK landing
% inside \textit -- which happens constantly, because \mtcitestyle
% italicises -- stays upright instead of drawing an "undefined shape".
\newcommand{\CJKface}[2]{%
  \newCJKfontfamily#1{NotoSansCJK-Regular.ttc}[%
    FontIndex = #2, ItalicFont = NotoSansCJK-Regular.ttc,
    ItalicFeatures = {FontIndex = #2}]}

\IfFontExistsTF{NotoSansCJK-Regular.ttc}{%
  \setCJKmainfont{NotoSansCJK-Regular.ttc}[%
    FontIndex = 2, ItalicFont = NotoSansCJK-Regular.ttc,
    ItalicFeatures = {FontIndex = 2}]
  \CJKface\chinesefontsf{2}	% Simplified
  \CJKface\yuefontsf{3}	% Traditional, used for Cantonese
  \CJKface\japanesefontsf{0}
  \CJKface\koreanfontsf{1}
}{%
  % No collection: fall back to the family names.  Everything still builds,
  % but the regional faces will all be whichever one fontconfig hands over.
  \setCJKmainfont{Noto Sans CJK SC}[ItalicFont = Noto Sans CJK SC]
  \newCJKfontfamily\chinesefontsf{Noto Sans CJK SC}[ItalicFont = Noto Sans CJK SC]
  \newCJKfontfamily\yuefontsf{Noto Sans CJK TC}[ItalicFont = Noto Sans CJK TC]
  \newCJKfontfamily\japanesefontsf{Noto Sans CJK JP}[ItalicFont = Noto Sans CJK JP]
  \newCJKfontfamily\koreanfontsf{Noto Sans CJK KR}[ItalicFont = Noto Sans CJK KR]
}
\setCJKsansfont{NotoSansCJK-Regular.ttc}[%
  FontIndex = 2, ItalicFont = NotoSansCJK-Regular.ttc,
  ItalicFeatures = {FontIndex = 2}]
% The mono faces live in the same collection: Mono JP=5, KR=6, SC=7, TC=8.
\setCJKmonofont{NotoSansCJK-Regular.ttc}[FontIndex = 7]

% The bopomofo tone marks sit in Spacing Modifier Letters, not the Bopomofo
% block, so xeCJK would leave them to Latin Modern -- which has no ˊ or ˋ.
\xeCJKDeclareCharClass{CJK}{"02C7, "02C9, "02CA, "02CB, "02D9}

% The reverse problem: xeCJK counts the curly quotes as CJK punctuation, so
% ' ' " " came out full-width in the CJK font, with CJK spacing around them
% and a line-break opportunity inside words like "doesn't".  In this course
% they are always English punctuation.
\xeCJKDeclareCharClass{Default}{"2018, "2019, "201C, "201D}

% IPA needs no font of its own any more: Biolinum covers IPA Extensions,
% Phonetic Extensions and the Spacing Modifier Letters.  The decks had grown
% four different answers to this (Andika, Noto Serif, Libertine, and nothing
% at all), so \ipa, \IPA and \silipa are all kept as names but now resolve
% to the body font.
\newfontfamily\ipafont{LibertinusSans}[
  Extension = .otf, UprightFont = *-Regular, ItalicFont = *-Italic,
  BoldFont = *-Bold, Ligatures = {TeX,NoCommon}]
\newcommand{\ipa}[1]{{\ipafont #1}}
\newcommand{\IPA}[1]{{\ipafont #1}}
\let\silipa\ipafont

% Used for the second line of gb4e glosses, where tone numbers and IPA turn
% up and Latin Modern runs out of superscripts.
\newfontfamily\Libertine{LibertinusSerif}[
  Extension = .otf, UprightFont = *-Regular, ItalicFont = *-Italic,
  BoldFont = *-Bold, BoldItalicFont = *-BoldItalic,
  Ligatures = {TeX,NoCommon}]

\usetheme{Madrid}
\usecolortheme{crane}

\usepackage{polyglossia}
\setdefaultlanguage{english}
\setotherlanguage{thai}
\newfontfamily\thaifont[Script=Thai]{Noto Sans Thai}
\newfontfamily\thaifontsf[Script=Thai]{Noto Sans Thai}
\setotherlanguage{arabic}
\newfontfamily\arabicfont[Script=Arabic]{Noto Naskh Arabic}
\newfontfamily\arabicfontsf[Script=Arabic]{Noto Naskh Arabic}

% polyglossia has no gloss file for these, so define the switches directly
% against the CJK families declared above.
\newcommand{\textchinese}[1]{{\chinesefontsf #1}}
\newcommand{\textjapanese}[1]{{\japanesefontsf #1}}
\newcommand{\textkorean}[1]{{\koreanfontsf #1}}
\newcommand{\textyue}[1]{{\yuefontsf #1}}

% Latin Modern has no mathematical angle brackets, which the writing systems
% deck uses for graphemes.
\usepackage{newunicodechar}
\newunicodechar{⟨}{‹}
\newunicodechar{⟩}{›}

%% ---- colour and emphasis ------------------------------------------------
\usepackage{xcolor}
\newcommand{\txx}[1]{\textcolor{blue}{\textbf{#1}}}
\let\emp\txx
\let\zz\txx
\newcommand{\mtplain}[1]{\textcolor{teal}{#1}}
\newcommand{\con}[1]{\textsc{#1}}

\usepackage[normalem]{ulem}
\newcommand{\ul}[1]{\uline{#1}}
\newcommand{\cll}[1]{\uline{#1}}

%% ---- linguistic examples ------------------------------------------------
%%
%% langsci-gb4e everywhere.  Plain gb4e makes _ and ^ active characters,
%% which breaks hyperref -- that is the one and only reason mygb4e existed, a
%% local copy of gb4e with exactly that hack commented out.  The Language
%% Science Press fork drops it, is maintained, and gives the \exfont /
%% \glossfont / \transfont hooks, so the local fork is no longer needed.
\usepackage{langsci-gb4e}

% Second gloss line -- transcription and IPA -- in the serif face.
% writing.tex and verbal-arts.tex override this afterwards.
\renewcommand{\eachwordtwo}{\Libertine}

%% ---- language shorthands ------------------------------------------------
%%
%% Two families, and the distinction matters:
%%
%%   ISO-code macros  \cmn{}, \jpn{} ... take a ROMANISATION and set it in
%%                    teal italic, optionally with a gloss: \cmn[dog]{gou3}.
%%   \zhs{} and friends take NATIVE SCRIPT and set it in teal upright in the
%%                    right regional face.  Script is never italicised.

\usepackage[c,e,f,j,k]{mtg2e}
\renewcommand{\mtcitestyle}[1]{\textcolor{teal}{\textit{#1}}}

% romanisations (mtg2e already gives \chn, \eng, \fra, \jpn, \kjpn and \kor)
\newcommand{\cmn}{\mtciteform}
\newcommand{\zh}{\mtciteform}
\newcommand{\yue}{\mtciteform}
\newcommand{\ko}{\mtciteform}
\newcommand{\vie}{\mtciteform}
\newcommand{\vi}{\mtciteform}
\newcommand{\tha}{\mtciteform}
\newcommand{\thaig}{\mtciteform}
\newcommand{\lao}{\mtciteform}
\newcommand{\ind}{\mtciteform}
\newcommand{\msa}{\mtciteform}
\newcommand{\zsm}{\mtciteform}
\newcommand{\ar}{\mtciteform}
\newcommand{\lex}[1]{\textbf{\mtcitestyle{#1}}}

% native script
\newcommand{\zhs}[1]{\mtplain{\textchinese{#1}}}
\newcommand{\zht}[1]{\mtplain{\textyue{#1}}}
\newcommand{\jps}[1]{\mtplain{\textjapanese{#1}}}
\newcommand{\kos}[1]{\mtplain{\textkorean{#1}}}
\newcommand{\ths}[1]{\mtplain{\textthai{#1}}}

%% ---- structure -----------------------------------------------------------
%% Sections are signposting, so give each one a page of its own showing where
%% we have got to, with the rest of the deck dimmed.  \overview puts the same
%% list up front; call it straight after \frame{\titlepage}.
\AtBeginSection[]{%
  \begin{frame}{Outline}
    \tableofcontents[currentsection]
  \end{frame}}

\newcommand{\overview}{%
  \begin{frame}{Overview}
    \tableofcontents
  \end{frame}}

%% ---- references ----------------------------------------------------------
\usepackage[round]{natbib}

%% ---- defaults ------------------------------------------------------------
%% The year lives here so the rollover is a one-line change.
\author[Francis Bond]{Francis Bond}
\date{2026}
 first, then
%% deviate, and put a comment saying why.

%% ---- fonts -------------------------------------------------------------
%%
%% Two mechanisms, deliberately, because the decks use both:
%%
%%   xeCJK      picks the CJK font automatically from the character, so Han,
%%              kana and Hangul just work in running text.
%%   polyglossia gives \textchinese{}, \textjapanese{}, \textkorean{},
%%              \textyue{}, \textthai{} and \textarabic{} for when a specific
%%              regional face is wanted -- the Japanese and Traditional faces
%%              draw many shared characters differently from the Simplified
%%              one, and Arabic needs RTL.
%%
%% Han has no italic.  Map the italic shape onto the upright one for every
%% CJK face, or XeTeX fake-slants it whenever CJK lands inside \textit --
%% which it constantly does, because \mtcitestyle italicises.

\usepackage{fontspec}
\usepackage{xeCJK}
\usepackage{graphicx}
\usepackage{tikz}

% Body font.  Latin Modern Sans, beamer's default, is the wrong font for this
% course: it has no small caps (so \textsc{sg} in a gloss came out as
% lowercase "sg") and it is missing 33 of the characters these decks use --
% most of the IPA, the modifier letters and the super/subscripts -- which
% XeTeX then drops silently.
%
% Libertinus is the maintained successor to Linux Libertine/Biolinum.  Sans
% becomes the body font (beamer sets \sffamily), Serif is there for gloss
% lines, and it brings a real maths font with it.  It lives in TeX Live, so
% kpathsea finds it on any machine without a fontconfig install.
% Set the faces directly rather than via libertinus-otf, which also wants a
% mono font (AnonymousPro) that is not installed here.
% Ligatures={TeX,NoCommon}: keep ``---'' and friends, but switch OFF the fi/fl
% output ligatures, so that text copied out of the PDF comes back as plain
% ASCII.  Those ligatures are how the deck sources picked up "Pacific" and
% "geneDcally" in the first place.
\usefonttheme{professionalfonts}
\setmainfont{LibertinusSerif}[
  Extension		= .otf,
  UprightFont		= *-Regular,
  ItalicFont		= *-Italic,
  BoldFont		= *-Bold,
  BoldItalicFont	= *-BoldItalic,
  Ligatures		= {TeX,NoCommon}]
% Scale 0.95: Libertinus has a much larger x-height than Latin Modern, so at
% the nominal size it overflows frames that used to fit.  0.95 restores every
% deck to the page count it had before the switch.
\setsansfont{LibertinusSans}[
  Scale			= 0.95,
  Extension		= .otf,
  UprightFont		= *-Regular,
  ItalicFont		= *-Italic,
  BoldFont		= *-Bold,
  % Libertinus Sans ships no bold italic, so slant the bold rather than let
  % LaTeX silently substitute an upright bold.
  BoldItalicFont	= *-Bold,
  BoldItalicFeatures	= {FakeSlant = 0.2},
  Ligatures		= {TeX,NoCommon}]
\setmonofont{LibertinusMono-Regular}[Extension = .otf]

% XeTeX always loads face 0 of a TrueType Collection, and
% NotoSansCJK-Regular.ttc packs JP=0, KR=1, SC=2, TC=3, HK=4.  So asking for
% "Noto Sans CJK SC" -- as every deck did -- silently handed back the
% JAPANESE face, and all the Simplified Chinese in this course was being set
% with Japanese glyph shapes (骨, 直, 关 and many more differ).  FontIndex
% picks the right face, and the bare filename resolves without a path.
%
% Han has no italic.  Point ItalicFont at the same face so that CJK landing
% inside \textit -- which happens constantly, because \mtcitestyle
% italicises -- stays upright instead of drawing an "undefined shape".
\newcommand{\CJKface}[2]{%
  \newCJKfontfamily#1{NotoSansCJK-Regular.ttc}[%
    FontIndex = #2, ItalicFont = NotoSansCJK-Regular.ttc,
    ItalicFeatures = {FontIndex = #2}]}

\IfFontExistsTF{NotoSansCJK-Regular.ttc}{%
  \setCJKmainfont{NotoSansCJK-Regular.ttc}[%
    FontIndex = 2, ItalicFont = NotoSansCJK-Regular.ttc,
    ItalicFeatures = {FontIndex = 2}]
  \CJKface\chinesefontsf{2}	% Simplified
  \CJKface\yuefontsf{3}	% Traditional, used for Cantonese
  \CJKface\japanesefontsf{0}
  \CJKface\koreanfontsf{1}
}{%
  % No collection: fall back to the family names.  Everything still builds,
  % but the regional faces will all be whichever one fontconfig hands over.
  \setCJKmainfont{Noto Sans CJK SC}[ItalicFont = Noto Sans CJK SC]
  \newCJKfontfamily\chinesefontsf{Noto Sans CJK SC}[ItalicFont = Noto Sans CJK SC]
  \newCJKfontfamily\yuefontsf{Noto Sans CJK TC}[ItalicFont = Noto Sans CJK TC]
  \newCJKfontfamily\japanesefontsf{Noto Sans CJK JP}[ItalicFont = Noto Sans CJK JP]
  \newCJKfontfamily\koreanfontsf{Noto Sans CJK KR}[ItalicFont = Noto Sans CJK KR]
}
\setCJKsansfont{NotoSansCJK-Regular.ttc}[%
  FontIndex = 2, ItalicFont = NotoSansCJK-Regular.ttc,
  ItalicFeatures = {FontIndex = 2}]
% The mono faces live in the same collection: Mono JP=5, KR=6, SC=7, TC=8.
\setCJKmonofont{NotoSansCJK-Regular.ttc}[FontIndex = 7]

% The bopomofo tone marks sit in Spacing Modifier Letters, not the Bopomofo
% block, so xeCJK would leave them to Latin Modern -- which has no ˊ or ˋ.
\xeCJKDeclareCharClass{CJK}{"02C7, "02C9, "02CA, "02CB, "02D9}

% The reverse problem: xeCJK counts the curly quotes as CJK punctuation, so
% ' ' " " came out full-width in the CJK font, with CJK spacing around them
% and a line-break opportunity inside words like "doesn't".  In this course
% they are always English punctuation.
\xeCJKDeclareCharClass{Default}{"2018, "2019, "201C, "201D}

% IPA needs no font of its own any more: Biolinum covers IPA Extensions,
% Phonetic Extensions and the Spacing Modifier Letters.  The decks had grown
% four different answers to this (Andika, Noto Serif, Libertine, and nothing
% at all), so \ipa, \IPA and \silipa are all kept as names but now resolve
% to the body font.
\newfontfamily\ipafont{LibertinusSans}[
  Extension = .otf, UprightFont = *-Regular, ItalicFont = *-Italic,
  BoldFont = *-Bold, Ligatures = {TeX,NoCommon}]
\newcommand{\ipa}[1]{{\ipafont #1}}
\newcommand{\IPA}[1]{{\ipafont #1}}
\let\silipa\ipafont

% Used for the second line of gb4e glosses, where tone numbers and IPA turn
% up and Latin Modern runs out of superscripts.
\newfontfamily\Libertine{LibertinusSerif}[
  Extension = .otf, UprightFont = *-Regular, ItalicFont = *-Italic,
  BoldFont = *-Bold, BoldItalicFont = *-BoldItalic,
  Ligatures = {TeX,NoCommon}]

\usetheme{Madrid}
\usecolortheme{crane}

\usepackage{polyglossia}
\setdefaultlanguage{english}
\setotherlanguage{thai}
\newfontfamily\thaifont[Script=Thai]{Noto Sans Thai}
\newfontfamily\thaifontsf[Script=Thai]{Noto Sans Thai}
\setotherlanguage{arabic}
\newfontfamily\arabicfont[Script=Arabic]{Noto Naskh Arabic}
\newfontfamily\arabicfontsf[Script=Arabic]{Noto Naskh Arabic}

% polyglossia has no gloss file for these, so define the switches directly
% against the CJK families declared above.
\newcommand{\textchinese}[1]{{\chinesefontsf #1}}
\newcommand{\textjapanese}[1]{{\japanesefontsf #1}}
\newcommand{\textkorean}[1]{{\koreanfontsf #1}}
\newcommand{\textyue}[1]{{\yuefontsf #1}}

% Latin Modern has no mathematical angle brackets, which the writing systems
% deck uses for graphemes.
\usepackage{newunicodechar}
\newunicodechar{⟨}{‹}
\newunicodechar{⟩}{›}

%% ---- colour and emphasis ------------------------------------------------
\usepackage{xcolor}
\newcommand{\txx}[1]{\textcolor{blue}{\textbf{#1}}}
\let\emp\txx
\let\zz\txx
\newcommand{\mtplain}[1]{\textcolor{teal}{#1}}
\newcommand{\con}[1]{\textsc{#1}}

\usepackage[normalem]{ulem}
\newcommand{\ul}[1]{\uline{#1}}
\newcommand{\cll}[1]{\uline{#1}}

%% ---- linguistic examples ------------------------------------------------
%%
%% langsci-gb4e everywhere.  Plain gb4e makes _ and ^ active characters,
%% which breaks hyperref -- that is the one and only reason mygb4e existed, a
%% local copy of gb4e with exactly that hack commented out.  The Language
%% Science Press fork drops it, is maintained, and gives the \exfont /
%% \glossfont / \transfont hooks, so the local fork is no longer needed.
\usepackage{langsci-gb4e}

% Second gloss line -- transcription and IPA -- in the serif face.
% writing.tex and verbal-arts.tex override this afterwards.
\renewcommand{\eachwordtwo}{\Libertine}

%% ---- language shorthands ------------------------------------------------
%%
%% Two families, and the distinction matters:
%%
%%   ISO-code macros  \cmn{}, \jpn{} ... take a ROMANISATION and set it in
%%                    teal italic, optionally with a gloss: \cmn[dog]{gou3}.
%%   \zhs{} and friends take NATIVE SCRIPT and set it in teal upright in the
%%                    right regional face.  Script is never italicised.

\usepackage[c,e,f,j,k]{mtg2e}
\renewcommand{\mtcitestyle}[1]{\textcolor{teal}{\textit{#1}}}

% romanisations (mtg2e already gives \chn, \eng, \fra, \jpn, \kjpn and \kor)
\newcommand{\cmn}{\mtciteform}
\newcommand{\zh}{\mtciteform}
\newcommand{\yue}{\mtciteform}
\newcommand{\ko}{\mtciteform}
\newcommand{\vie}{\mtciteform}
\newcommand{\vi}{\mtciteform}
\newcommand{\tha}{\mtciteform}
\newcommand{\thaig}{\mtciteform}
\newcommand{\lao}{\mtciteform}
\newcommand{\ind}{\mtciteform}
\newcommand{\msa}{\mtciteform}
\newcommand{\zsm}{\mtciteform}
\newcommand{\ar}{\mtciteform}
\newcommand{\lex}[1]{\textbf{\mtcitestyle{#1}}}

% native script
\newcommand{\zhs}[1]{\mtplain{\textchinese{#1}}}
\newcommand{\zht}[1]{\mtplain{\textyue{#1}}}
\newcommand{\jps}[1]{\mtplain{\textjapanese{#1}}}
\newcommand{\kos}[1]{\mtplain{\textkorean{#1}}}
\newcommand{\ths}[1]{\mtplain{\textthai{#1}}}

%% ---- structure -----------------------------------------------------------
%% Sections are signposting, so give each one a page of its own showing where
%% we have got to, with the rest of the deck dimmed.  \overview puts the same
%% list up front; call it straight after \frame{\titlepage}.
\AtBeginSection[]{%
  \begin{frame}{Outline}
    \tableofcontents[currentsection]
  \end{frame}}

\newcommand{\overview}{%
  \begin{frame}{Overview}
    \tableofcontents
  \end{frame}}

%% ---- references ----------------------------------------------------------
\usepackage[round]{natbib}

%% ---- defaults ------------------------------------------------------------
%% The year lives here so the rollover is a one-line change.
\author[Francis Bond]{Francis Bond}
\date{2026}
 first, then
%% deviate, and put a comment saying why.

%% ---- fonts -------------------------------------------------------------
%%
%% Two mechanisms, deliberately, because the decks use both:
%%
%%   xeCJK      picks the CJK font automatically from the character, so Han,
%%              kana and Hangul just work in running text.
%%   polyglossia gives \textchinese{}, \textjapanese{}, \textkorean{},
%%              \textyue{}, \textthai{} and \textarabic{} for when a specific
%%              regional face is wanted -- the Japanese and Traditional faces
%%              draw many shared characters differently from the Simplified
%%              one, and Arabic needs RTL.
%%
%% Han has no italic.  Map the italic shape onto the upright one for every
%% CJK face, or XeTeX fake-slants it whenever CJK lands inside \textit --
%% which it constantly does, because \mtcitestyle italicises.

\usepackage{fontspec}
\usepackage{xeCJK}
\usepackage{graphicx}
\usepackage{tikz}

% Body font.  Latin Modern Sans, beamer's default, is the wrong font for this
% course: it has no small caps (so \textsc{sg} in a gloss came out as
% lowercase "sg") and it is missing 33 of the characters these decks use --
% most of the IPA, the modifier letters and the super/subscripts -- which
% XeTeX then drops silently.
%
% Libertinus is the maintained successor to Linux Libertine/Biolinum.  Sans
% becomes the body font (beamer sets \sffamily), Serif is there for gloss
% lines, and it brings a real maths font with it.  It lives in TeX Live, so
% kpathsea finds it on any machine without a fontconfig install.
% Set the faces directly rather than via libertinus-otf, which also wants a
% mono font (AnonymousPro) that is not installed here.
% Ligatures={TeX,NoCommon}: keep ``---'' and friends, but switch OFF the fi/fl
% output ligatures, so that text copied out of the PDF comes back as plain
% ASCII.  Those ligatures are how the deck sources picked up "Pacific" and
% "geneDcally" in the first place.
\usefonttheme{professionalfonts}
\setmainfont{LibertinusSerif}[
  Extension		= .otf,
  UprightFont		= *-Regular,
  ItalicFont		= *-Italic,
  BoldFont		= *-Bold,
  BoldItalicFont	= *-BoldItalic,
  Ligatures		= {TeX,NoCommon}]
% Scale 0.95: Libertinus has a much larger x-height than Latin Modern, so at
% the nominal size it overflows frames that used to fit.  0.95 restores every
% deck to the page count it had before the switch.
\setsansfont{LibertinusSans}[
  Scale			= 0.95,
  Extension		= .otf,
  UprightFont		= *-Regular,
  ItalicFont		= *-Italic,
  BoldFont		= *-Bold,
  % Libertinus Sans ships no bold italic, so slant the bold rather than let
  % LaTeX silently substitute an upright bold.
  BoldItalicFont	= *-Bold,
  BoldItalicFeatures	= {FakeSlant = 0.2},
  Ligatures		= {TeX,NoCommon}]
\setmonofont{LibertinusMono-Regular}[Extension = .otf]

% XeTeX always loads face 0 of a TrueType Collection, and
% NotoSansCJK-Regular.ttc packs JP=0, KR=1, SC=2, TC=3, HK=4.  So asking for
% "Noto Sans CJK SC" -- as every deck did -- silently handed back the
% JAPANESE face, and all the Simplified Chinese in this course was being set
% with Japanese glyph shapes (骨, 直, 关 and many more differ).  FontIndex
% picks the right face, and the bare filename resolves without a path.
%
% Han has no italic.  Point ItalicFont at the same face so that CJK landing
% inside \textit -- which happens constantly, because \mtcitestyle
% italicises -- stays upright instead of drawing an "undefined shape".
\newcommand{\CJKface}[2]{%
  \newCJKfontfamily#1{NotoSansCJK-Regular.ttc}[%
    FontIndex = #2, ItalicFont = NotoSansCJK-Regular.ttc,
    ItalicFeatures = {FontIndex = #2}]}

\IfFontExistsTF{NotoSansCJK-Regular.ttc}{%
  \setCJKmainfont{NotoSansCJK-Regular.ttc}[%
    FontIndex = 2, ItalicFont = NotoSansCJK-Regular.ttc,
    ItalicFeatures = {FontIndex = 2}]
  \CJKface\chinesefontsf{2}	% Simplified
  \CJKface\yuefontsf{3}	% Traditional, used for Cantonese
  \CJKface\japanesefontsf{0}
  \CJKface\koreanfontsf{1}
}{%
  % No collection: fall back to the family names.  Everything still builds,
  % but the regional faces will all be whichever one fontconfig hands over.
  \setCJKmainfont{Noto Sans CJK SC}[ItalicFont = Noto Sans CJK SC]
  \newCJKfontfamily\chinesefontsf{Noto Sans CJK SC}[ItalicFont = Noto Sans CJK SC]
  \newCJKfontfamily\yuefontsf{Noto Sans CJK TC}[ItalicFont = Noto Sans CJK TC]
  \newCJKfontfamily\japanesefontsf{Noto Sans CJK JP}[ItalicFont = Noto Sans CJK JP]
  \newCJKfontfamily\koreanfontsf{Noto Sans CJK KR}[ItalicFont = Noto Sans CJK KR]
}
\setCJKsansfont{NotoSansCJK-Regular.ttc}[%
  FontIndex = 2, ItalicFont = NotoSansCJK-Regular.ttc,
  ItalicFeatures = {FontIndex = 2}]
% The mono faces live in the same collection: Mono JP=5, KR=6, SC=7, TC=8.
\setCJKmonofont{NotoSansCJK-Regular.ttc}[FontIndex = 7]

% The bopomofo tone marks sit in Spacing Modifier Letters, not the Bopomofo
% block, so xeCJK would leave them to Latin Modern -- which has no ˊ or ˋ.
\xeCJKDeclareCharClass{CJK}{"02C7, "02C9, "02CA, "02CB, "02D9}

% The reverse problem: xeCJK counts the curly quotes as CJK punctuation, so
% ' ' " " came out full-width in the CJK font, with CJK spacing around them
% and a line-break opportunity inside words like "doesn't".  In this course
% they are always English punctuation.
\xeCJKDeclareCharClass{Default}{"2018, "2019, "201C, "201D}

% IPA needs no font of its own any more: Biolinum covers IPA Extensions,
% Phonetic Extensions and the Spacing Modifier Letters.  The decks had grown
% four different answers to this (Andika, Noto Serif, Libertine, and nothing
% at all), so \ipa, \IPA and \silipa are all kept as names but now resolve
% to the body font.
\newfontfamily\ipafont{LibertinusSans}[
  Extension = .otf, UprightFont = *-Regular, ItalicFont = *-Italic,
  BoldFont = *-Bold, Ligatures = {TeX,NoCommon}]
\newcommand{\ipa}[1]{{\ipafont #1}}
\newcommand{\IPA}[1]{{\ipafont #1}}
\let\silipa\ipafont

% Used for the second line of gb4e glosses, where tone numbers and IPA turn
% up and Latin Modern runs out of superscripts.
\newfontfamily\Libertine{LibertinusSerif}[
  Extension = .otf, UprightFont = *-Regular, ItalicFont = *-Italic,
  BoldFont = *-Bold, BoldItalicFont = *-BoldItalic,
  Ligatures = {TeX,NoCommon}]

\usetheme{Madrid}
\usecolortheme{crane}

\usepackage{polyglossia}
\setdefaultlanguage{english}
\setotherlanguage{thai}
\newfontfamily\thaifont[Script=Thai]{Noto Sans Thai}
\newfontfamily\thaifontsf[Script=Thai]{Noto Sans Thai}
\setotherlanguage{arabic}
\newfontfamily\arabicfont[Script=Arabic]{Noto Naskh Arabic}
\newfontfamily\arabicfontsf[Script=Arabic]{Noto Naskh Arabic}

% polyglossia has no gloss file for these, so define the switches directly
% against the CJK families declared above.
\newcommand{\textchinese}[1]{{\chinesefontsf #1}}
\newcommand{\textjapanese}[1]{{\japanesefontsf #1}}
\newcommand{\textkorean}[1]{{\koreanfontsf #1}}
\newcommand{\textyue}[1]{{\yuefontsf #1}}

% Latin Modern has no mathematical angle brackets, which the writing systems
% deck uses for graphemes.
\usepackage{newunicodechar}
\newunicodechar{⟨}{‹}
\newunicodechar{⟩}{›}

%% ---- colour and emphasis ------------------------------------------------
\usepackage{xcolor}
\newcommand{\txx}[1]{\textcolor{blue}{\textbf{#1}}}
\let\emp\txx
\let\zz\txx
\newcommand{\mtplain}[1]{\textcolor{teal}{#1}}
\newcommand{\con}[1]{\textsc{#1}}

\usepackage[normalem]{ulem}
\newcommand{\ul}[1]{\uline{#1}}
\newcommand{\cll}[1]{\uline{#1}}

%% ---- linguistic examples ------------------------------------------------
%%
%% langsci-gb4e everywhere.  Plain gb4e makes _ and ^ active characters,
%% which breaks hyperref -- that is the one and only reason mygb4e existed, a
%% local copy of gb4e with exactly that hack commented out.  The Language
%% Science Press fork drops it, is maintained, and gives the \exfont /
%% \glossfont / \transfont hooks, so the local fork is no longer needed.
\usepackage{langsci-gb4e}

% Second gloss line -- transcription and IPA -- in the serif face.
% writing.tex and verbal-arts.tex override this afterwards.
\renewcommand{\eachwordtwo}{\Libertine}

%% ---- language shorthands ------------------------------------------------
%%
%% Two families, and the distinction matters:
%%
%%   ISO-code macros  \cmn{}, \jpn{} ... take a ROMANISATION and set it in
%%                    teal italic, optionally with a gloss: \cmn[dog]{gou3}.
%%   \zhs{} and friends take NATIVE SCRIPT and set it in teal upright in the
%%                    right regional face.  Script is never italicised.

\usepackage[c,e,f,j,k]{mtg2e}
\renewcommand{\mtcitestyle}[1]{\textcolor{teal}{\textit{#1}}}

% romanisations (mtg2e already gives \chn, \eng, \fra, \jpn, \kjpn and \kor)
\newcommand{\cmn}{\mtciteform}
\newcommand{\zh}{\mtciteform}
\newcommand{\yue}{\mtciteform}
\newcommand{\ko}{\mtciteform}
\newcommand{\vie}{\mtciteform}
\newcommand{\vi}{\mtciteform}
\newcommand{\tha}{\mtciteform}
\newcommand{\thaig}{\mtciteform}
\newcommand{\lao}{\mtciteform}
\newcommand{\ind}{\mtciteform}
\newcommand{\msa}{\mtciteform}
\newcommand{\zsm}{\mtciteform}
\newcommand{\ar}{\mtciteform}
\newcommand{\lex}[1]{\textbf{\mtcitestyle{#1}}}

% native script
\newcommand{\zhs}[1]{\mtplain{\textchinese{#1}}}
\newcommand{\zht}[1]{\mtplain{\textyue{#1}}}
\newcommand{\jps}[1]{\mtplain{\textjapanese{#1}}}
\newcommand{\kos}[1]{\mtplain{\textkorean{#1}}}
\newcommand{\ths}[1]{\mtplain{\textthai{#1}}}

%% ---- structure -----------------------------------------------------------
%% Sections are signposting, so give each one a page of its own showing where
%% we have got to, with the rest of the deck dimmed.  \overview puts the same
%% list up front; call it straight after \frame{\titlepage}.
\AtBeginSection[]{%
  \begin{frame}{Outline}
    \tableofcontents[currentsection]
  \end{frame}}

\newcommand{\overview}{%
  \begin{frame}{Overview}
    \tableofcontents
  \end{frame}}

%% ---- references ----------------------------------------------------------
\usepackage[round]{natbib}

%% ---- defaults ------------------------------------------------------------
%% The year lives here so the rollover is a one-line change.
\author[Francis Bond]{Francis Bond}
\date{2026}
 first, then
%% deviate, and put a comment saying why.

%% ---- fonts -------------------------------------------------------------
%%
%% Two mechanisms, deliberately, because the decks use both:
%%
%%   xeCJK      picks the CJK font automatically from the character, so Han,
%%              kana and Hangul just work in running text.
%%   polyglossia gives \textchinese{}, \textjapanese{}, \textkorean{},
%%              \textyue{}, \textthai{} and \textarabic{} for when a specific
%%              regional face is wanted -- the Japanese and Traditional faces
%%              draw many shared characters differently from the Simplified
%%              one, and Arabic needs RTL.
%%
%% Han has no italic.  Map the italic shape onto the upright one for every
%% CJK face, or XeTeX fake-slants it whenever CJK lands inside \textit --
%% which it constantly does, because \mtcitestyle italicises.

\usepackage{fontspec}
\usepackage{xeCJK}
\usepackage{graphicx}
\usepackage{tikz}

% Body font.  Latin Modern Sans, beamer's default, is the wrong font for this
% course: it has no small caps (so \textsc{sg} in a gloss came out as
% lowercase "sg") and it is missing 33 of the characters these decks use --
% most of the IPA, the modifier letters and the super/subscripts -- which
% XeTeX then drops silently.
%
% Libertinus is the maintained successor to Linux Libertine/Biolinum.  Sans
% becomes the body font (beamer sets \sffamily), Serif is there for gloss
% lines, and it brings a real maths font with it.  It lives in TeX Live, so
% kpathsea finds it on any machine without a fontconfig install.
% Set the faces directly rather than via libertinus-otf, which also wants a
% mono font (AnonymousPro) that is not installed here.
% Ligatures={TeX,NoCommon}: keep ``---'' and friends, but switch OFF the fi/fl
% output ligatures, so that text copied out of the PDF comes back as plain
% ASCII.  Those ligatures are how the deck sources picked up "Pacific" and
% "geneDcally" in the first place.
\usefonttheme{professionalfonts}
\setmainfont{LibertinusSerif}[
  Extension		= .otf,
  UprightFont		= *-Regular,
  ItalicFont		= *-Italic,
  BoldFont		= *-Bold,
  BoldItalicFont	= *-BoldItalic,
  Ligatures		= {TeX,NoCommon}]
% Scale 0.95: Libertinus has a much larger x-height than Latin Modern, so at
% the nominal size it overflows frames that used to fit.  0.95 restores every
% deck to the page count it had before the switch.
\setsansfont{LibertinusSans}[
  Scale			= 0.95,
  Extension		= .otf,
  UprightFont		= *-Regular,
  ItalicFont		= *-Italic,
  BoldFont		= *-Bold,
  % Libertinus Sans ships no bold italic, so slant the bold rather than let
  % LaTeX silently substitute an upright bold.
  BoldItalicFont	= *-Bold,
  BoldItalicFeatures	= {FakeSlant = 0.2},
  Ligatures		= {TeX,NoCommon}]
\setmonofont{LibertinusMono-Regular}[Extension = .otf]

% XeTeX always loads face 0 of a TrueType Collection, and
% NotoSansCJK-Regular.ttc packs JP=0, KR=1, SC=2, TC=3, HK=4.  So asking for
% "Noto Sans CJK SC" -- as every deck did -- silently handed back the
% JAPANESE face, and all the Simplified Chinese in this course was being set
% with Japanese glyph shapes (骨, 直, 关 and many more differ).  FontIndex
% picks the right face, and the bare filename resolves without a path.
%
% Han has no italic.  Point ItalicFont at the same face so that CJK landing
% inside \textit -- which happens constantly, because \mtcitestyle
% italicises -- stays upright instead of drawing an "undefined shape".
\newcommand{\CJKface}[2]{%
  \newCJKfontfamily#1{NotoSansCJK-Regular.ttc}[%
    FontIndex = #2, ItalicFont = NotoSansCJK-Regular.ttc,
    ItalicFeatures = {FontIndex = #2}]}

\IfFontExistsTF{NotoSansCJK-Regular.ttc}{%
  \setCJKmainfont{NotoSansCJK-Regular.ttc}[%
    FontIndex = 2, ItalicFont = NotoSansCJK-Regular.ttc,
    ItalicFeatures = {FontIndex = 2}]
  \CJKface\chinesefontsf{2}	% Simplified
  \CJKface\yuefontsf{3}	% Traditional, used for Cantonese
  \CJKface\japanesefontsf{0}
  \CJKface\koreanfontsf{1}
}{%
  % No collection: fall back to the family names.  Everything still builds,
  % but the regional faces will all be whichever one fontconfig hands over.
  \setCJKmainfont{Noto Sans CJK SC}[ItalicFont = Noto Sans CJK SC]
  \newCJKfontfamily\chinesefontsf{Noto Sans CJK SC}[ItalicFont = Noto Sans CJK SC]
  \newCJKfontfamily\yuefontsf{Noto Sans CJK TC}[ItalicFont = Noto Sans CJK TC]
  \newCJKfontfamily\japanesefontsf{Noto Sans CJK JP}[ItalicFont = Noto Sans CJK JP]
  \newCJKfontfamily\koreanfontsf{Noto Sans CJK KR}[ItalicFont = Noto Sans CJK KR]
}
\setCJKsansfont{NotoSansCJK-Regular.ttc}[%
  FontIndex = 2, ItalicFont = NotoSansCJK-Regular.ttc,
  ItalicFeatures = {FontIndex = 2}]
% The mono faces live in the same collection: Mono JP=5, KR=6, SC=7, TC=8.
\setCJKmonofont{NotoSansCJK-Regular.ttc}[FontIndex = 7]

% The bopomofo tone marks sit in Spacing Modifier Letters, not the Bopomofo
% block, so xeCJK would leave them to Latin Modern -- which has no ˊ or ˋ.
\xeCJKDeclareCharClass{CJK}{"02C7, "02C9, "02CA, "02CB, "02D9}

% The reverse problem: xeCJK counts the curly quotes as CJK punctuation, so
% ' ' " " came out full-width in the CJK font, with CJK spacing around them
% and a line-break opportunity inside words like "doesn't".  In this course
% they are always English punctuation.
\xeCJKDeclareCharClass{Default}{"2018, "2019, "201C, "201D}

% IPA needs no font of its own any more: Biolinum covers IPA Extensions,
% Phonetic Extensions and the Spacing Modifier Letters.  The decks had grown
% four different answers to this (Andika, Noto Serif, Libertine, and nothing
% at all), so \ipa, \IPA and \silipa are all kept as names but now resolve
% to the body font.
\newfontfamily\ipafont{LibertinusSans}[
  Extension = .otf, UprightFont = *-Regular, ItalicFont = *-Italic,
  BoldFont = *-Bold, Ligatures = {TeX,NoCommon}]
\newcommand{\ipa}[1]{{\ipafont #1}}
\newcommand{\IPA}[1]{{\ipafont #1}}
\let\silipa\ipafont

% Used for the second line of gb4e glosses, where tone numbers and IPA turn
% up and Latin Modern runs out of superscripts.
\newfontfamily\Libertine{LibertinusSerif}[
  Extension = .otf, UprightFont = *-Regular, ItalicFont = *-Italic,
  BoldFont = *-Bold, BoldItalicFont = *-BoldItalic,
  Ligatures = {TeX,NoCommon}]

\usetheme{Madrid}
\usecolortheme{crane}

\usepackage{polyglossia}
\setdefaultlanguage{english}
\setotherlanguage{thai}
\newfontfamily\thaifont[Script=Thai]{Noto Sans Thai}
\newfontfamily\thaifontsf[Script=Thai]{Noto Sans Thai}
\setotherlanguage{arabic}
\newfontfamily\arabicfont[Script=Arabic]{Noto Naskh Arabic}
\newfontfamily\arabicfontsf[Script=Arabic]{Noto Naskh Arabic}

% polyglossia has no gloss file for these, so define the switches directly
% against the CJK families declared above.
\newcommand{\textchinese}[1]{{\chinesefontsf #1}}
\newcommand{\textjapanese}[1]{{\japanesefontsf #1}}
\newcommand{\textkorean}[1]{{\koreanfontsf #1}}
\newcommand{\textyue}[1]{{\yuefontsf #1}}

% Latin Modern has no mathematical angle brackets, which the writing systems
% deck uses for graphemes.
\usepackage{newunicodechar}
\newunicodechar{⟨}{‹}
\newunicodechar{⟩}{›}

%% ---- colour and emphasis ------------------------------------------------
\usepackage{xcolor}
\newcommand{\txx}[1]{\textcolor{blue}{\textbf{#1}}}
\let\emp\txx
\let\zz\txx
\newcommand{\mtplain}[1]{\textcolor{teal}{#1}}
\newcommand{\con}[1]{\textsc{#1}}

\usepackage[normalem]{ulem}
\newcommand{\ul}[1]{\uline{#1}}
\newcommand{\cll}[1]{\uline{#1}}

%% ---- linguistic examples ------------------------------------------------
%%
%% langsci-gb4e everywhere.  Plain gb4e makes _ and ^ active characters,
%% which breaks hyperref -- that is the one and only reason mygb4e existed, a
%% local copy of gb4e with exactly that hack commented out.  The Language
%% Science Press fork drops it, is maintained, and gives the \exfont /
%% \glossfont / \transfont hooks, so the local fork is no longer needed.
\usepackage{langsci-gb4e}

% Second gloss line -- transcription and IPA -- in the serif face.
% writing.tex and verbal-arts.tex override this afterwards.
\renewcommand{\eachwordtwo}{\Libertine}

%% ---- language shorthands ------------------------------------------------
%%
%% Two families, and the distinction matters:
%%
%%   ISO-code macros  \cmn{}, \jpn{} ... take a ROMANISATION and set it in
%%                    teal italic, optionally with a gloss: \cmn[dog]{gou3}.
%%   \zhs{} and friends take NATIVE SCRIPT and set it in teal upright in the
%%                    right regional face.  Script is never italicised.

\usepackage[c,e,f,j,k]{mtg2e}
\renewcommand{\mtcitestyle}[1]{\textcolor{teal}{\textit{#1}}}

% romanisations (mtg2e already gives \chn, \eng, \fra, \jpn, \kjpn and \kor)
\newcommand{\cmn}{\mtciteform}
\newcommand{\zh}{\mtciteform}
\newcommand{\yue}{\mtciteform}
\newcommand{\ko}{\mtciteform}
\newcommand{\vie}{\mtciteform}
\newcommand{\vi}{\mtciteform}
\newcommand{\tha}{\mtciteform}
\newcommand{\thaig}{\mtciteform}
\newcommand{\lao}{\mtciteform}
\newcommand{\ind}{\mtciteform}
\newcommand{\msa}{\mtciteform}
\newcommand{\zsm}{\mtciteform}
\newcommand{\ar}{\mtciteform}
\newcommand{\lex}[1]{\textbf{\mtcitestyle{#1}}}

% native script
\newcommand{\zhs}[1]{\mtplain{\textchinese{#1}}}
\newcommand{\zht}[1]{\mtplain{\textyue{#1}}}
\newcommand{\jps}[1]{\mtplain{\textjapanese{#1}}}
\newcommand{\kos}[1]{\mtplain{\textkorean{#1}}}
\newcommand{\ths}[1]{\mtplain{\textthai{#1}}}

%% ---- structure -----------------------------------------------------------
%% Sections are signposting, so give each one a page of its own showing where
%% we have got to, with the rest of the deck dimmed.  \overview puts the same
%% list up front; call it straight after \frame{\titlepage}.
\AtBeginSection[]{%
  \begin{frame}{Outline}
    \tableofcontents[currentsection]
  \end{frame}}

\newcommand{\overview}{%
  \begin{frame}{Overview}
    \tableofcontents
  \end{frame}}

%% ---- references ----------------------------------------------------------
\usepackage[round]{natbib}

%% ---- defaults ------------------------------------------------------------
%% The year lives here so the rollover is a one-line change.
\author[Francis Bond]{Francis Bond}
\date{2026}
